% Options for packages loaded elsewhere
% Options for packages loaded elsewhere
\PassOptionsToPackage{unicode}{hyperref}
\PassOptionsToPackage{hyphens}{url}
\PassOptionsToPackage{dvipsnames,svgnames,x11names}{xcolor}
%
\documentclass[
  letterpaper,
  DIV=11,
  numbers=noendperiod]{scrartcl}
\usepackage{xcolor}
\usepackage{amsmath,amssymb}
\setcounter{secnumdepth}{5}
\usepackage{iftex}
\ifPDFTeX
  \usepackage[T1]{fontenc}
  \usepackage[utf8]{inputenc}
  \usepackage{textcomp} % provide euro and other symbols
\else % if luatex or xetex
  \usepackage{unicode-math} % this also loads fontspec
  \defaultfontfeatures{Scale=MatchLowercase}
  \defaultfontfeatures[\rmfamily]{Ligatures=TeX,Scale=1}
\fi
\usepackage{lmodern}
\ifPDFTeX\else
  % xetex/luatex font selection
\fi
% Use upquote if available, for straight quotes in verbatim environments
\IfFileExists{upquote.sty}{\usepackage{upquote}}{}
\IfFileExists{microtype.sty}{% use microtype if available
  \usepackage[]{microtype}
  \UseMicrotypeSet[protrusion]{basicmath} % disable protrusion for tt fonts
}{}
\makeatletter
\@ifundefined{KOMAClassName}{% if non-KOMA class
  \IfFileExists{parskip.sty}{%
    \usepackage{parskip}
  }{% else
    \setlength{\parindent}{0pt}
    \setlength{\parskip}{6pt plus 2pt minus 1pt}}
}{% if KOMA class
  \KOMAoptions{parskip=half}}
\makeatother
% Make \paragraph and \subparagraph free-standing
\makeatletter
\ifx\paragraph\undefined\else
  \let\oldparagraph\paragraph
  \renewcommand{\paragraph}{
    \@ifstar
      \xxxParagraphStar
      \xxxParagraphNoStar
  }
  \newcommand{\xxxParagraphStar}[1]{\oldparagraph*{#1}\mbox{}}
  \newcommand{\xxxParagraphNoStar}[1]{\oldparagraph{#1}\mbox{}}
\fi
\ifx\subparagraph\undefined\else
  \let\oldsubparagraph\subparagraph
  \renewcommand{\subparagraph}{
    \@ifstar
      \xxxSubParagraphStar
      \xxxSubParagraphNoStar
  }
  \newcommand{\xxxSubParagraphStar}[1]{\oldsubparagraph*{#1}\mbox{}}
  \newcommand{\xxxSubParagraphNoStar}[1]{\oldsubparagraph{#1}\mbox{}}
\fi
\makeatother


\usepackage{longtable,booktabs,array}
\newcounter{none} % for unnumbered tables
\usepackage{calc} % for calculating minipage widths
% Correct order of tables after \paragraph or \subparagraph
\usepackage{etoolbox}
\makeatletter
\patchcmd\longtable{\par}{\if@noskipsec\mbox{}\fi\par}{}{}
\makeatother
% Allow footnotes in longtable head/foot
\IfFileExists{footnotehyper.sty}{\usepackage{footnotehyper}}{\usepackage{footnote}}
\makesavenoteenv{longtable}
\usepackage{graphicx}
\makeatletter
\newsavebox\pandoc@box
\newcommand*\pandocbounded[1]{% scales image to fit in text height/width
  \sbox\pandoc@box{#1}%
  \Gscale@div\@tempa{\textheight}{\dimexpr\ht\pandoc@box+\dp\pandoc@box\relax}%
  \Gscale@div\@tempb{\linewidth}{\wd\pandoc@box}%
  \ifdim\@tempb\p@<\@tempa\p@\let\@tempa\@tempb\fi% select the smaller of both
  \ifdim\@tempa\p@<\p@\scalebox{\@tempa}{\usebox\pandoc@box}%
  \else\usebox{\pandoc@box}%
  \fi%
}
% Set default figure placement to htbp
\def\fps@figure{htbp}
\makeatother


% definitions for citeproc citations
\NewDocumentCommand\citeproctext{}{}
\NewDocumentCommand\citeproc{mm}{%
  \begingroup\def\citeproctext{#2}\cite{#1}\endgroup}
\makeatletter
 % allow citations to break across lines
 \let\@cite@ofmt\@firstofone
 % avoid brackets around text for \cite:
 \def\@biblabel#1{}
 \def\@cite#1#2{{#1\if@tempswa , #2\fi}}
\makeatother
\newlength{\cslhangindent}
\setlength{\cslhangindent}{1.5em}
\newlength{\csllabelwidth}
\setlength{\csllabelwidth}{3em}
\newenvironment{CSLReferences}[2] % #1 hanging-indent, #2 entry-spacing
 {\begin{list}{}{%
  \setlength{\itemindent}{0pt}
  \setlength{\leftmargin}{0pt}
  \setlength{\parsep}{0pt}
  % turn on hanging indent if param 1 is 1
  \ifodd #1
   \setlength{\leftmargin}{\cslhangindent}
   \setlength{\itemindent}{-1\cslhangindent}
  \fi
  % set entry spacing
  \setlength{\itemsep}{#2\baselineskip}}}
 {\end{list}}
\usepackage{calc}
\newcommand{\CSLBlock}[1]{\hfill\break\parbox[t]{\linewidth}{\strut\ignorespaces#1\strut}}
\newcommand{\CSLLeftMargin}[1]{\parbox[t]{\csllabelwidth}{\strut#1\strut}}
\newcommand{\CSLRightInline}[1]{\parbox[t]{\linewidth - \csllabelwidth}{\strut#1\strut}}
\newcommand{\CSLIndent}[1]{\hspace{\cslhangindent}#1}



\setlength{\emergencystretch}{3em} % prevent overfull lines

\providecommand{\tightlist}{%
  \setlength{\itemsep}{0pt}\setlength{\parskip}{0pt}}



 


\usepackage{booktabs}
\usepackage{longtable}
\usepackage{array}
\usepackage{multirow}
\usepackage{wrapfig}
\usepackage{float}
\usepackage{colortbl}
\usepackage{pdflscape}
\usepackage{tabu}
\usepackage{threeparttable}
\usepackage{threeparttablex}
\usepackage[normalem]{ulem}
\usepackage{makecell}
\usepackage{xcolor}
\usepackage{tabularray}
\usepackage[normalem]{ulem}
\usepackage{graphicx}
\usepackage{rotating}
\UseTblrLibrary{siunitx}
\NewTableCommand{\tinytableDefineColor}[3]{\definecolor{#1}{#2}{#3}}
\newcommand{\tinytableTabularrayUnderline}[1]{\underline{#1}}
\newcommand{\tinytableTabularrayStrikeout}[1]{\sout{#1}}
\usepackage{siunitx}

    \newcolumntype{d}{S[
      table-align-text-before=false,
      table-align-text-after=false,
      input-symbols={-,\*+()}
    ]}
  
\KOMAoption{captions}{tableheading}
\usepackage[font=small,labelfont=bf]{caption}
\sisetup{detect-all}
\raggedbottom
\interfootnotelinepenalty=10000
\setlength{\skip\footins}{2\baselineskip}
\makeatletter
\@ifpackageloaded{caption}{}{\usepackage{caption}}
\AtBeginDocument{%
\ifdefined\contentsname
  \renewcommand*\contentsname{Table of contents}
\else
  \newcommand\contentsname{Table of contents}
\fi
\ifdefined\listfigurename
  \renewcommand*\listfigurename{List of Figures}
\else
  \newcommand\listfigurename{List of Figures}
\fi
\ifdefined\listtablename
  \renewcommand*\listtablename{List of Tables}
\else
  \newcommand\listtablename{List of Tables}
\fi
\ifdefined\figurename
  \renewcommand*\figurename{Figure}
\else
  \newcommand\figurename{Figure}
\fi
\ifdefined\tablename
  \renewcommand*\tablename{Table}
\else
  \newcommand\tablename{Table}
\fi
}
\@ifpackageloaded{float}{}{\usepackage{float}}
\floatstyle{ruled}
\@ifundefined{c@chapter}{\newfloat{codelisting}{h}{lop}}{\newfloat{codelisting}{h}{lop}[chapter]}
\floatname{codelisting}{Listing}
\newcommand*\listoflistings{\listof{codelisting}{List of Listings}}
\makeatother
\makeatletter
\makeatother
\makeatletter
\@ifpackageloaded{caption}{}{\usepackage{caption}}
\@ifpackageloaded{subcaption}{}{\usepackage{subcaption}}
\makeatother
\usepackage{bookmark}
\IfFileExists{xurl.sty}{\usepackage{xurl}}{} % add URL line breaks if available
\urlstyle{same}
\hypersetup{
  pdftitle={Characterizing Canadian Out-of-district Political Donations from 2015 to 2024},
  pdfauthor={Benedict Cummins-Mburu; Rohan Alexander},
  colorlinks=true,
  linkcolor={blue},
  filecolor={Maroon},
  citecolor={Blue},
  urlcolor={Blue},
  pdfcreator={LaTeX via pandoc}}


\title{Characterizing Canadian Out-of-district Political Donations from
2015 to 2024\thanks{Code and data are available at:
\url{https://zenodo.org/records/21781460}. We thank Natural Sciences and
Engineering Research Council of Canada (NSERC) Alliance Grant (ALLRP
599949 - 24) and Social Sciences and Humanities Research Council (SSHRC)
Partnership Grant (895-2025-1002) for financial support. We thank Chris
Cochrane for suggesting this area of research, and Martin Allen with his
help procuring our dataset. We also thank Ellie Murray, Inessa De
Angelis, Mariana Garcia Mejia, Oscar Heath, Sabrina Kreyzerman, and Zane
Schwartz for their helpful comments. Corresponding author:
\href{mailto:rohan.alexander@utoronto.ca}{\nolinkurl{rohan.alexander@utoronto.ca}}.}}
\author{Benedict Cummins-Mburu \and Rohan Alexander}
\date{August 3, 2026}
\begin{document}
\maketitle
\begin{abstract}
Federal political representation in Canada is geographical: one Member
of Parliament (MP) represents one electoral district. Despite this
representational structure, Canadians can legally donate to a host of
political entities, regardless of proximity. Using a public dataset from
Elections Canada's Open Data on Contributions matched with census data
from Statistics Canada, we characterize broad patterns in Canadian
out-of-district donation from 2015 to 2024. We find that out-of-district
donations form a persistent share of donor-sourced funding to both
candidates and district associations, averaging 40\% across time,
election periods, and party lines. While relatively constant through
time, the propensity for out-of-district donation varies substantially
across districts, beyond what spatial proximity alone would predict.
Districts with higher population density, median household income,
educational attainment, and visible minority proportions see higher
rates of out-of-district giving, the same traits that predict donor
prevalence generally in Canada. Despite its frequency, out-of-district
donation in Canada remains largely regionalized; nearly half of
donations that leave their district are received by a political entity
in a neighbouring one, and few cross provincial lines. These and other
findings provide a helpful characterization of this largely undocumented
facet of Canadian political financing, and inform further research on
effective representation in Canada and other nations with similar
representational structures.\\
\textbf{Keywords:} out-of-district donation; Canadian politics;
political geography; political financing
\end{abstract}


\newpage

\section{Introduction}\label{sec-intro}

In Canada, federal representation is structurally geographical; citizens
are grouped into a few hundred electoral districts based on where they
live, each represented by a single Member of Parliament (MP). An
individual's political influence in these systems is canonically
exercised by voting for their parliamentary candidate of choice during
election periods, however alternate forms of engagement exist to support
a broader scope of political activities, political entities, and without
the spatio-temporal constraints imposed by elections. Citizens may
volunteer at their local district associations, advocate for their
preferred national party online, vote in intra-party nomination
contests, or contribute financially (i.e.~donate) to a candidate's
campaign (\citeproc{ref-bednar_political_2011}{Bednar and Gerber 2011};
\citeproc{ref-baker_getting_2016}{Baker 2016};
\citeproc{ref-garnett_where_2024}{Garnett 2024}).

Whether intentional or not, these alternative forms of political
participation may serve to extend a citizen's political influence beyond
their district's boundaries. Among these, out-of-district donations are
perhaps the clearest and most flexible mechanism through which this is
achieved, for two main reasons. First, donations are largely
unconstrained in time and space. In North America, citizens and
permanent residents can legally contribute to various federal political
entities at any time, in any district, and as often as they want,
subject to annual and per-election limitations
(\citeproc{ref-gimpel_check_2008}{Gimpel et al. 2008};
\citeproc{ref-canada_annual_2026}{Elections Canada 2026a}). Second,
individual contributions are a major source of funding for electoral
campaigns (\citeproc{ref-gimpel_check_2008}{Gimpel et al. 2008};
\citeproc{ref-garnett_where_2024}{Garnett 2024};
\citeproc{ref-canada_understanding_2026}{Elections Canada 2026f}), and a
major source of revenue for national parties
(\citeproc{ref-corrado_new_2010}{Corrado et al. 2010};
\citeproc{ref-garnett_lifeblood_2022}{Garnett et al. 2022}). In the
U.S., even incumbent representatives often attribute a majority of their
income to their individual donor base
(\citeproc{ref-baker_getting_2016}{Baker 2016}). While the activities of
these political actors may be reimbursed through public stipends and
subsidies, the gross funding needed for, say, a party leadership
contestant to carry out an effective campaign, often needs to be
available on-demand (\citeproc{ref-garnett_where_2024}{Garnett 2024}).
Despite a growing body of Canadian research documenting the extent of
and patterns of federal political donations generally (e.g.
\citeproc{ref-besco_ethnic_2022}{Besco and Tolley 2022};
\citeproc{ref-tolley_who_2022}{Tolley et al. 2022};
\citeproc{ref-garnett_where_2024}{Garnett 2024}), comparatively little
attention has been given to out-of-district donations.

Existing research documents substantial rates of out-of-district federal
contributions in the U.S. Both congressional and senatorial candidates
presently source between 50 and 70 percent of their campaign funding
from outside their district, a proportion that has increased in recent
decades (\citeproc{ref-gimpel_check_2008}{Gimpel et al. 2008};
\citeproc{ref-bednar_political_2011}{Bednar and Gerber 2011};
\citeproc{ref-sievert_out_2023}{Sievert and Mathiasen 2023}). Incumbent
representatives also receive around half of their donations from
out-of-district, which meaningfully influences their roll-call behaviour
(\citeproc{ref-peoples_impact_2008}{Peoples and Gortari 2008};
\citeproc{ref-canes_wrone_out_2022}{Canes-Wrone and Miller 2022}). These
donations are thought to be primarily motivated by partisan gain,
evidenced, for example, by increased donations to competitive districts
(\citeproc{ref-gimpel_check_2008}{Gimpel et al. 2008};
\citeproc{ref-jacobs_nationalization_2024}{Jacobs and Imboywa 2024}).
Funding is also geographically concentrated: both major parties in the
U.S. primarily source local funding from a single small group of highly
influential and socio-demographically similar districts
(\citeproc{ref-gimpel_political_2006}{Gimpel et al. 2006}).

In this paper, we document the extent of, and large-scale trends in,
Canadian out-of-district donations to electoral candidates and district
associations from 2015 to 2024. We use a database containing all
contributions received by federal political entities during this period,
made publicly accessible by \emph{Elections Canada}. This dataset was
augmented by situating donors geographically, and matching donations to
their socio-demographic environment using the 2021 Census, allowing us
to better understand which factors correlate with out-of-district
donation rates.

We find that, between 2015 and 2024, out-of-district donations account
for a persistent share of donor-sourced funding to both electoral
candidates and district associations, averaging roughly 40\% across
time, election periods, and party lines. This relative invariance does
not extend to space, however: the propensity for out-of-district
donation changes substantially from district to district. Donations
originating from densely populated, high-income, and well-educated
districts are more likely to leave their district of origin, echoing
predictors of out-of-district giving established in the U.S. Our
analysis also shows that the visible minority proportion of a donor's
district is the most correlated socio-demographic variable that we
considered. This finding departs from the ethnic profile historically
associated with ``donor districts'' in the U.S.
(\citeproc{ref-gimpel_check_2008}{Gimpel et al. 2008};
\citeproc{ref-canes_wrone_out_2022}{Canes-Wrone and Miller 2022}), but
aligns with recent Canadian research linking minority ethnic communities
to higher rates of political participation and out-of-district giving
(\citeproc{ref-garnett_lifeblood_2022}{Garnett et al. 2022};
\citeproc{ref-besco_ethnic_2022}{Besco and Tolley 2022}). Despite the
considerable presence of out-of-district donations during our study
period, out-of-district giving remains a regionalized phenomenon: nearly
half of all donations that leave their district of origin are addressed
to political entities in a neighbouring one, and fewer still cross a
provincial boundary, limiting their potential influence on
geographically-based representation. Although we do not formally
evaluate the representational consequences of out-of-district donations,
the systematic characterization presented here provides a foundation for
future research on political finance and effective representation in
Canada, while also advancing understanding of how political influence is
distributed across the electorate.

\section{Background}\label{sec-lit}

\subsection{Motivation: Representational Effects of Out-of-district
Donations in the U.S.}\label{sec-lit-motivation}

The extent, trends, and socio-demographic characteristics of U.S.
out-of-district donations are well- documented among individual
contributions to campaigning and incumbent House representatives (both
senators and congressional representatives). This body of research was
catalyzed by Gimpel et al. (\citeproc{ref-gimpel_check_2008}{2008})'s
examination of out-of-district donations during four Congressional
elections between 1994 and 2004, and their impacts on legislative
representation. They propose the existence of ``donor districts''; a
small set of older, wealthy, well-educated, and urban districts wherein
a large majority of out-of-district donations originate, destined to
Republican and Democrat candidates alike. Notably, these
socio-demographic predictors of out-of-district donor output are largely
the same as predictors of donor output generally in both the U.S. and
Canada (\citeproc{ref-gimpel_political_2006}{Gimpel et al. 2006};
\citeproc{ref-garnett_where_2024}{Garnett 2024}). Over this time period,
they also observed an increase in the proportion of non-local donations
received by the average candidate, from around 50\% in 1994 to nearly
70\% in 2004 (\citeproc{ref-gimpel_check_2008}{Gimpel et al. 2008}).
More recent studies suggest that this and similar trends have continued
past 2004 (\citeproc{ref-bramlett_political_2011}{Bramlett et al. 2011};
\citeproc{ref-baker_getting_2016}{Baker 2016};
\citeproc{ref-canes_wrone_out_2022}{Canes-Wrone and Miller 2022}), and
for donations made during U.S. senate elections
(\citeproc{ref-sievert_out_2023}{Sievert and Mathiasen 2023}) Driven by
the observation that the ``closeness'' of local races strongly predicted
which districts received these non-local donations, Gimpel et al.
(\citeproc{ref-gimpel_check_2008}{2008}) argue that gaining a partisan
advantage (and not seeking ideological or identity-based representation)
is the primary motivation behind non-local donations.

Follow-up U.S. research has found that out-of-district campaign
contributions affect the voting behaviour of Congressional
representatives, thereby distorting the supposedly-dyadic relationship
between a representative and their district. A study of the 2006-2010
Congressional elections revealed that the average congressional
representative becomes less responsive to their constituencies'
ideological preferences with increased out-of-district contributions
(\citeproc{ref-baker_getting_2016}{Baker 2016}). Canes-Wrone and Miller
(\citeproc{ref-canes_wrone_out_2022}{2022}) further explore the
2006-2016 Congressional elections to show that House voting behaviour is
indeed associated with the preferences of their national donor base,
contrasted against research finding that political action committees do
not affect roll call behaviour
(\citeproc{ref-ansolabehere_why_2003}{Ansolabehere et al. 2003}).
Another study on senatorial elections emphasizes the rivalrous
relationship between constituents and out-of-district donors,
experimentally showing that constituents who learned about outside
donation expected their senators to be less loyal to local interests
(\citeproc{ref-nathan_donor_2025}{Nathan et al. 2025}). More generally,
longitudinal analyses of out-of-state (and therefore out-of-district)
donations during senatorial elections echo many of the findings from
congressional elections (\citeproc{ref-sievert_out_2023}{Sievert and
Mathiasen 2023}; \citeproc{ref-jacobs_nationalization_2024}{Jacobs and
Imboywa 2024}; \citeproc{ref-keena_out_2024}{Keena 2024}).

Strikingly, out-of-district federal political donations have not been
analyzed to nearly the same extent in other Westminster-style
parliamentary democracies, including Canada. Besides differences in
population size and the scale of electoral competition, we propose two
barriers to research in the Canadian context. First, U.S. annual
contribution limits are much higher than Canadian ones, but the
threshold of location disclosure is the same
(\citeproc{ref-gimpel_political_2006}{Gimpel et al. 2006};
\citeproc{ref-fec_contribution_limits_2026}{Federal Election Commission
2026}; \citeproc{ref-canada_understanding_2026}{Elections Canada
2026f}). This means that location-disclosed donations (and therefore
those that can be analyzed as out-of-district) likely constitute a
larger share of the total amount of all donor-based federal political
financing. Second, the data manipulation steps required to situate
donors within districts and attribute accurate Census-based data to the
entries are more straightforward in the U.S. as compared to Canada. U.S.
Census information is available directly at the zip-code level, and
their geographic conversion files (to situate zip-codes within
districts) are fully open-source. While donor geolocation and
Census-based data augmentation has been done in some Canadian research
(e.g. \citeproc{ref-garnett_where_2024}{Garnett 2024}), most of the
Canadian donation studies we identified rely on political surveys such
as the Canadian Elections Study.

\subsection{Brief Review of Donations to Canadian Electoral
Candidates}\label{sec-lit-review}

Limited accounts of out-of-district donation rates exist in the Canadian
context. One study analyzing co-ethnic affinity among donors and
candidates found that the likelihood of a donor donating out-of-district
is higher for ethnic minorities (50-70\%) compared to donors of European
origin (33\%) (\citeproc{ref-besco_ethnic_2022}{Besco and Tolley 2022}).
However, these proportions only describe candidate-bound donations
during the 2015 general election period. Here we review existing
research concerning donations to Canadian electoral candidates,
contextualizing our subsequent analysis\footnote{Although donations to
  district associations are also considered in our study, we were unable
  to find research focusing on contributions addressed to these.}. We do
not consider studies focusing on donations to other political entities
as these are outside the scope of our study\footnote{We refer the reader
  to Young and Jansen (\citeproc{ref-young_money_2011}{2011}), Garnett
  et al. (\citeproc{ref-garnett_lifeblood_2022}{2022}), and Pruysers
  (\citeproc{ref-pruysers_leadership_2024}{2024}) for examples of such
  studies.}.

Considerable research has examined socio-demographic characteristics of
donors and candidates in Canada. Preliminary analyses of individual
donations from 2006 to 2011 (\citeproc{ref-jansen_donates_2012}{Jansen
et al. 2012}; \citeproc{ref-carmichael_political_2014}{Carmichael and
Howe 2014}), and a more recent analysis of candidate-bound donations
during the 2015, 2019, and 2021 general elections
(\citeproc{ref-garnett_where_2024}{Garnett 2024}) have found that donors
are more likely to originate from older, wealthier, and better-educated
districts. The gender gap in giving has also received attention, likely
given the broader context of established gender gaps in political
participation (\citeproc{ref-tolley_who_2022}{Tolley et al. 2022};
\citeproc{ref-mcmahon_covid_2023}{McMahon et al. 2023}). From 2004 to
2015, overall federal donation amounts have increased, particularly in
Western Canada, and among Conservative political entities
(\citeproc{ref-mcmahon_follow_2015}{McMahon and Wolfe 2015}). However,
contribution rates to candidates during the following three general
elections (in 2015, 2019, and 2021) have since decreased
(\citeproc{ref-garnett_where_2024}{Garnett 2024}), and the COVID-19
pandemic saw overall reduced political participation
(\citeproc{ref-mcmahon_covid_2023}{McMahon et al. 2023}).

Canadian donors tend to donate more to candidates that they share
aspects of their identity with. For instance, women donors are more
likely to donate to women candidates
(\citeproc{ref-tolley_who_2022}{Tolley et al. 2022}), and there are high
rates of co-ethnic donation, particularly among South Asian donors
(\citeproc{ref-besco_ethnic_2022}{Besco and Tolley 2022}). More
generally, South Asian political participation appears to be
particularly heightened in Canada, relative to other ethnocultural
groups (\citeproc{ref-garnett_where_2024}{Garnett 2024}).

Some studies have also attempted to assess the impacts of political
donations on representation. Among incumbent legislators, voting
behaviour is seemingly independent of the composition of their donor
base (\citeproc{ref-peoples_impact_2008}{Peoples and Gortari 2008}).
However, Eagles (\citeproc{ref-eagles_effectiveness_2004}{2004}) found
that local campaign donations during the 1993 and 1997 elections
meaningfully affected the number of votes a candidate receives.

\section{Data and Methods}\label{sec-data}

\subsection{Data Overview}\label{sec-data-overview}

The data used in this paper was originally sourced from the Elections
Canada Open Data on Contributions datasets
(\citeproc{ref-canada_open_2025}{Elections Canada 2025a}). Since 2004,
the \emph{Canada Elections Act} has required that all contributions
(monetary or otherwise) valued at over \$20 be recorded by the recipient
and shared with the government
(\citeproc{ref-canada_understanding_2026}{Elections Canada 2026f}). The
Elections Canada dataset on donations contains all donations made by
individuals to registered federal political entities since 2004
(\citeproc{ref-canada_open_2025}{Elections Canada 2025a}). Additionally,
donors who contributed over \$200 must have their full name and address
(including their postal code) recorded
(\citeproc{ref-canada_understanding_2026}{Elections Canada 2026f}). For
all contributions, Elections Canada also records the recipient's name,
affiliated district (if applicable), party (if applicable), entity
(e.g.~national party, candidate, etc.), and the date the contribution
was received. More information on political financing laws in Canada is
presented in Appendix~\ref{sec-info}. The primary dataset used in this
study is a cleaned version of Elections Canada's, made available to us
by the Investigative Journalism Foundation. Their cleaning process is
described in Appendix~\ref{sec-pipeline}.

To even be considered as out-of-district, a federal donation must be
received by a political entity that is representationally tied to a
specific district. For this reason, only donations received by
candidates during election periods and district associations\footnote{Although
  nomination contestants and by-election candidates are also local
  political entities and were present in the data, these were not
  considered to simplify our analysis and discussion, since they only
  make up 0.7\% of donations addressed to localized entities.}
(hereafter referred to as localized entities) can be considered in the
study of out-of-district donations. Donations received by non-localized
entities present in this dataset (parties or leadership contestants)
cannot. Given their prevalence, however, we briefly compare these with
localized entities in Section~\ref{sec-data-summaries-situation}.

Every ten years, federal electoral district (FED) boundaries are redrawn
to account for population changes that occurred within the decade
(\citeproc{ref-canada_redistribution_2025}{Elections Canada 2025b}). To
restrict our focus to a single, static set of FEDs, our analysis only
includes contributions made during the 2013 representational order,
effective between October 2, 2015 and April 22, 2024
(\citeproc{ref-canada_redistribution_2025}{Elections Canada 2025b}).
This is the largest such order present in the dataset, capturing 61\% of
all recorded contributions since 2004, and spanning three general
election periods (42nd, 43rd, and 44th). Donors were matched to their
FEDs by their postal codes using Statistics Canada's Postal Code
Conversion Files\footnote{Postal codes may straddle multiple FEDs. For
  this study, we created a deterministic spatial lookup table that
  resolves these cases using reasonable criteria. Donations whose postal
  codes could not be matched in this way were dropped. Our
  methodologies, and statistical assessments of their impacts are
  detailed further in Appendix~\ref{sec-data-supplements-spatial}.}
(\citeproc{ref-statcan_pccf_2025}{Statistics Canada 2025}), made
available to us by the University of Toronto. Donor FEDs were then
compared with recipient FEDs to assess whether a donation was
out-of-district. Donor location was not recorded for the majority of
donations under \$200, and the small share that could be
donor-geolocated were distributed very differently from the rest of the
data. For these reasons, this study only considers donations above
\$200, as was done in similar studies
(\citeproc{ref-gimpel_check_2008}{Gimpel et al. 2008};
\citeproc{ref-garnett_where_2024}{Garnett 2024}). These smaller
donations are discussed further in Appendix~\ref{sec-missing-data}.

Donations were also situated within census dissemination areas (DA),
allowing us to attribute socio-demographic characteristics to donation
sources at a fine spatial grain. For each dissemination area, we
identified its population, population density, median household income,
mean age, proportion of the population with a post-secondary education,
and proportions of various visible minority ethnicities. A full list of
the ethnicities is available in Appendix~\ref{sec-data-supplements}.
These variables were chosen since they are associated with the
propensity to donate during an election
(\citeproc{ref-gimpel_political_2006}{Gimpel et al. 2006},
\citeproc{ref-gimpel_check_2008}{2008};
\citeproc{ref-garnett_where_2024}{Garnett 2024}). Further, ethnicities
were disaggregated since previous Canadian research has found that some
minority groups such as South Asians donate at much higher rates than
others (\citeproc{ref-besco_ethnic_2022}{Besco and Tolley 2022}). We
again used the PCCF to match donor postal codes to the DAs of the 2021
Census, and to group DAs into FEDs\footnote{Dissemination areas may
  straddle multiple FEDs. Appendix~\ref{sec-data-supplements-spatial}
  discusses how these cases were handled.} to deduce the same
socio-demographic information at the district level. A detailed
walkthrough of the full data pipeline, from the files available for
download on Canadian government websites to the core datasets used for
analysis can be found in Appendix~\ref{sec-pipeline} and
Appendix~\ref{sec-data-supplements}. All data cleaning, visualizations,
and analyses were conducted using the R statistical software
(\citeproc{ref-r_citation}{R Core Team 2025}).

\subsection{Data Summaries}\label{sec-data-summaries}

\subsubsection{Situating Localized
Donations}\label{sec-data-summaries-situation}

We first provide preliminary summaries of all contributions\footnote{Throughout
  this paper, the terms ``contribution'' and ``donation'' are used
  interchangeably and implicitly refer to contributions above \$200 that
  were successfully donor-geolocated, unless otherwise specified.}
during our study period, to help situate further analysis of
out-of-district donations. Table~\ref{tbl-data-summaries-situation}
breaks down donations by recipient type. National entities (parties and
leadership contestants) received the majority (80.8\%) of contributions
during our study period. Among donations addressed to localized
entities, 24.3\% are received by candidates, and 75.7\% were received by
district associations (EDAs)\footnote{In our dataset, candidates receive
  all of their donations during writ-periods. During these periods, they
  instead represent 48\% of the total amount contributed to localized
  entities, with district associations making up the remaining 52\%
  (Appendix~\ref{sec-additional-data-summaries})}. Proportions are
similar when broken down by count instead of amount.

Figure~\ref{fig-data-summaries-situation-timelines} shows how
contributions during our study period are distributed in
time\footnote{All donations to electoral candidates and leadership
  contestants were tied to specific events in time. A small proportion
  (0.7\% of the total contributed amount) of contributions to district
  associations and national parties over \$200 were missing specific
  dates, but could always be attributed to a fiscal year. Plots over
  time in this report include these cases, which slightly exaggerates
  the observed December peaks in contribution rates.}, for national and
localized entities separately. Among national-level donations, we
observe three distinct temporal signatures. First, donation rates are
higher in December of every year during our study, likely since
charitable donations tend to be more common in this period in general,
as individuals claim tax deductions and give more during the holidays.
Second, we observe six peaks corresponding to general election periods
(peaks during and immediately before grey shaded regions), and
Conservative party leadership campaigns (peaks in the blue trend line in
the first half of 2017, 2020, and 2022). Third, national-entity donation
rates generally increased over our study period, following the
increasing trend observed from 2004 to 2015 by McMahon and Wolfe
(\citeproc{ref-mcmahon_follow_2015}{2015}) during the previous decade.
This is particularly noticeable after the 2021 general election, and for
donations received by Conservative national entities.

Donation rates to localized entities exhibited the same December peaks,
and peaks during general election periods that were observed for
donations to national entities. However, localized donation rates did
not exhibit an increasing trend, unlike their counterparts bound for
national entities. We also observe a significant reduction in donation
rates from the start of 2019 to mid-2021, which may be partially
explained by the COVID-19 pandemic, agreeing with previous findings on
reduced contribution rates during this period
(\citeproc{ref-mcmahon_covid_2023}{McMahon et al. 2023}). Spatial
comparisons between the aggregated donations bound for national
vs.~localized entities during the study period yielded only minor
differences, and are presented in
Appendix~\ref{sec-additional-data-summaries}, which did not yield
drastic differences.

\begin{table}

\caption{\label{tbl-data-summaries-situation}Summaries of financial
contributions to federal political entities over \$200 from 2015 to
2024, separated by entity. Only localized contributions (those addressed
to candidates and district associations) are considered further in our
study of out-of-district donations. Contributions received by nomination
contestants are not shown, nor studied.}

\centering{

\centering
\begin{tblr}[         %% tabularray outer open
]                     %% tabularray outer close
{                     %% tabularray inner open
width={1\linewidth},
colspec={X[]X[]X[]X[]X[]},
hline{2}={1-5}{solid, black, 0.05em},
hline{1}={1-5}{solid, black, 0.08em},
hline{6}={1-5}{solid, black, 0.08em},
cell{1}{1}={}{font=\bfseries\fontsize{0.85em}{1.15em}\selectfont, halign=l},
cell{1}{2}={}{font=\bfseries\fontsize{0.85em}{1.15em}\selectfont, halign=l},
cell{1}{3}={}{font=\bfseries\fontsize{0.85em}{1.15em}\selectfont, halign=r},
cell{1}{4}={}{font=\bfseries\fontsize{0.85em}{1.15em}\selectfont, halign=r},
cell{1}{5}={}{font=\bfseries\fontsize{0.85em}{1.15em}\selectfont, halign=r},
cell{2-5}{1}={}{font=\fontsize{0.85em}{1.15em}\selectfont, halign=l},
cell{2-5}{2}={}{font=\fontsize{0.85em}{1.15em}\selectfont, halign=l},
cell{2-5}{3}={}{font=\fontsize{0.85em}{1.15em}\selectfont, halign=r},
cell{2-5}{4}={}{font=\fontsize{0.85em}{1.15em}\selectfont, halign=r},
cell{2-5}{5}={}{font=\fontsize{0.85em}{1.15em}\selectfont, halign=r},
}                     %% tabularray inner close
Political Entity & Type & Donation Count & Donation Amount (Thousands) & Share of Total Amount \\
Candidate & Localized & 30,559 & \$19,883 & 4.6\% \\
District Association & Localized & 108,619 & \$63,354 & 14.6\% \\
Leadership Contestant & National & 47,873 & \$32,451 & 7.5\% \\
Party & National & 618,609 & \$319,056 & 73.4\% \\
\end{tblr}

}

\end{table}%

\begin{figure}

\centering{

\pandocbounded{\includegraphics[keepaspectratio]{paper_files/figure-pdf/fig-data-summaries-situation-timelines-1.pdf}}

}

\caption{\label{fig-data-summaries-situation-timelines}Timeline of
financial contributions to federal political entities over \$200, from
2015 to 2024. (A) shows only contributions to national entities (parties
and leadership contestants), and (B) shows only contributions to
localized entities (candidates and district associations). Within each
panel, trendlines are separated by the party of the recipient. Shaded
vertical bars indicate the 42nd, 43rd, and 44th general election periods
respectively. Graphs were qualitatively similar when contributions under
\$200 were included.}

\end{figure}%

\subsubsection{Spatio-temporal Trends in Out-of-District Donation
Rates}\label{sec-data-summaries-ood-rate}

In this section, we present data summaries describing when and where
out-of-district contributions are more or less common.
Table~\ref{tbl-data-summaries-ood-rate-overall} displays the proportion
of contributions above \$200 (raw and amount-weighted) that were sent
out-of-district, for candidate- and EDA-bound donations. Overall,
candidates received 37.6\% of these donations from outside their
district, and EDAs received 43.8\%\footnote{A comparison of the value
  distributions of in- and out-of-district donations is available in
  Appendix~\ref{sec-additional-data-summaries}.}. In both cases, these
proportions increase when donations are weighed by their value,
indicating that out-of-district donations are larger, on average, than
their within-district counterparts.

Table~\ref{tbl-data-summaries-ood-rate-candidates} and
Figure~\ref{fig-data-summaries-ood-rate-EDAs} show how out-of-district
funding rates change over time, for candidates and district associations
respectively\footnote{It is unclear why donations to EDAs plummet in
  2019 in our dataset. This does not affect our analysis however, since
  we never use dates directly.}. Candidate out-of-district rates show a
slight decreasing trend, from 39.5\% in the 42nd general election in
2015 to around 35.4\% in subsequent periods
(Table~\ref{tbl-data-summaries-ood-rate-candidates}). However, this may
reflect the heightened absolute number of contributions garnered in 2015
compared to other periods. EDAs received roughly the same proportion of
out-of-district contributions during the entire study period, including
during writ periods\footnote{Roughly similar rates, and the same
  consistency through time, were observed when entities from individual
  parties were examined (Appendix~\ref{sec-additional-data-summaries}).}.

\begin{table}

\caption{\label{tbl-data-summaries-ood-rate-overall}Proportion of
contributions over \$200 that were received by localized entities
(candidates and district associations) outside the donor's district,
from 2015 to 2024. Raw and value-weighted proportions of out-of-district
donations (OOD) are reported. Overall contribution counts and amount are
also reported for reference. Contributions received by candidates
outside general election cycles (by-elections) are not included in
calculations.}

\centering{

\centering
\begin{tblr}[         %% tabularray outer open
]                     %% tabularray outer close
{                     %% tabularray inner open
width={1\linewidth},
colspec={X[]X[]X[]X[]X[]},
hline{2}={1-5}{solid, black, 0.05em},
hline{1}={1-5}{solid, black, 0.08em},
hline{4}={1-5}{solid, black, 0.08em},
cell{1}{1}={}{font=\bfseries\fontsize{0.85em}{1.15em}\selectfont, halign=l},
cell{1}{2}={}{font=\bfseries\fontsize{0.85em}{1.15em}\selectfont, halign=l},
cell{1}{3}={}{font=\bfseries\fontsize{0.85em}{1.15em}\selectfont, halign=r},
cell{1}{4}={}{font=\bfseries\fontsize{0.85em}{1.15em}\selectfont, halign=r},
cell{1}{5}={}{font=\bfseries\fontsize{0.85em}{1.15em}\selectfont, halign=r},
cell{2-3}{1}={}{font=\fontsize{0.85em}{1.15em}\selectfont, halign=l},
cell{2-3}{2}={}{font=\fontsize{0.85em}{1.15em}\selectfont, halign=l},
cell{2-3}{3}={}{font=\fontsize{0.85em}{1.15em}\selectfont, halign=r},
cell{2-3}{4}={}{font=\fontsize{0.85em}{1.15em}\selectfont, halign=r},
cell{2-3}{5}={}{font=\fontsize{0.85em}{1.15em}\selectfont, halign=r},
}                     %% tabularray inner close
Political Entity & Donation Count & Donation Amount (Thousands) & OOD Proportion & OOD Proportion (Amount Weighted) \\
Candidate & 30,559 & \$19,883 & 37.6\% & 40.2\% \\
District Association & 108,619 & \$63,354 & 43.8\% & 47.8\% \\
\end{tblr}

}

\end{table}%

\begin{table}

\caption{\label{tbl-data-summaries-ood-rate-candidates}Proportion of
contributions to electoral candidates over \$200 that were
out-of-district (OOD) over the three general election periods that
occurred during our study period. Raw and value-weighted proportions are
reported, along with the overall number of contributions (in- and
out-of-district) for reference. Contributions received by candidates
outside general election cycles (by-elections) are not included in
calculations.}

\centering{

\centering
\begin{tblr}[         %% tabularray outer open
]                     %% tabularray outer close
{                     %% tabularray inner open
width={1\linewidth},
colspec={X[]X[]X[]X[]X[]},
hline{2}={1-5}{solid, black, 0.05em},
hline{1}={1-5}{solid, black, 0.08em},
hline{5}={1-5}{solid, black, 0.08em},
cell{1}{1}={}{font=\bfseries\fontsize{0.85em}{1.15em}\selectfont, halign=l},
cell{1}{2}={}{font=\bfseries\fontsize{0.85em}{1.15em}\selectfont, halign=l},
cell{1}{3}={}{font=\bfseries\fontsize{0.85em}{1.15em}\selectfont, halign=r},
cell{1}{4}={}{font=\bfseries\fontsize{0.85em}{1.15em}\selectfont, halign=r},
cell{1}{5}={}{font=\bfseries\fontsize{0.85em}{1.15em}\selectfont, halign=r},
cell{2-4}{1}={}{font=\fontsize{0.85em}{1.15em}\selectfont, halign=l},
cell{2-4}{2}={}{font=\fontsize{0.85em}{1.15em}\selectfont, halign=l},
cell{2-4}{3}={}{font=\fontsize{0.85em}{1.15em}\selectfont, halign=r},
cell{2-4}{4}={}{font=\fontsize{0.85em}{1.15em}\selectfont, halign=r},
cell{2-4}{5}={}{font=\fontsize{0.85em}{1.15em}\selectfont, halign=r},
}                     %% tabularray inner close
General Election & Donation Count & Donation Amount (Thousands) & OOD Proportion & OOD Proportion (Amount Weighted) \\
42nd general election & 16,662 & \$10,383 & 39.5\% & 41.9\% \\
43rd general election & 7,900 & \$5,204 & 35.2\% & 39.1\% \\
44th general election & 5,997 & \$4,296 & 35.6\% & 37.7\% \\
\end{tblr}

}

\end{table}%

\begin{figure}

\centering{

\pandocbounded{\includegraphics[keepaspectratio]{paper_files/figure-pdf/fig-data-summaries-ood-rate-EDAs-1.pdf}}

}

\caption{\label{fig-data-summaries-ood-rate-EDAs}Timeline of the
out-of-district share of the total monthly contributed amount over \$200
addressed to district associations from 2015 to 2024. Raw monthly
proportions are displayed as grey points, and the fitted trend over time
(with its 95\% confidence interval) is shown with the blue line and
ribbon. Trend over time was fit using a generalized additive model with
a smooth term over donation date, estimated at the individual level.
This is shown instead of a simple line plot to account for the large
month-to-month variation in donation rates to district associations
generally. Shaded vertical bars indicate the 42nd, 43rd, and 44th
general election periods respectively. The horizontal dashed line
represents the overall share of contributed amount that was received
from out-of-district during the study period, for reference.}

\end{figure}%

Figure~\ref{fig-data-summaries-ood-rate-district-scatterplot} displays
the relationship between in- and out-of-district contribution amounts,
for each electoral district present in our dataset. We display donations
to candidates and district associations separately. In both cases,
out-of-district contribution rates vary greatly from district to
district\footnote{The marginal distribution of out-of-district
  proportions across all districts is shown in
  Appendix~\ref{sec-additional-data-summaries}.}. Some Toronto
metropolitan districts such as Toronto---St.~Paul's and
University---Rosedale sent between 70\% and 80\% of their contributions
to localized entities in other districts, while others such as Alberta's
Lethbridge and Ontario's Niagara Falls only sent between 5\% and 15\% of
their localized contribution sum to other districts. Many districts,
especially for donations addressed to candidates, have low
out-of-district donation rates, as illustrated by the line of points
near \(y = 0\) in
Figure~\ref{fig-data-summaries-ood-rate-district-scatterplot}. We also
find that, at least among the few top districts by contribution amount,
out-of-district propensity is not a simple function of
urbanization\footnote{Since districts vary enormously in size and shape
  depending on their population density, one might naturally expect that
  tightly-packed urban districts would have higher out-of-district
  contribution rates, simply due to the spatial nature of
  socio-political interest. In
  Appendix~\ref{sec-district-propensity-simulation}, we simulate the
  out-of-district share for each district under this expectation, and
  find that district-level out-of-district proportions differ much more
  than what can be explained by the spatial arrangement of districts
  alone.}; Vancouver Granville in British Columbia and Brampton North
are similarly metropolitan, but have very different out-of-district
contribution rates to electoral candidates\footnote{District-level
  differences in candidate-bound out-of-district contribution rates are
  even more pronounced when election periods are examined individually
  (Appendix~\ref{sec-additional-data-summaries}).}.

\begin{figure}

\centering{

\pandocbounded{\includegraphics[keepaspectratio]{paper_files/figure-pdf/fig-data-summaries-ood-rate-district-scatterplot-1.pdf}}

}

\caption{\label{fig-data-summaries-ood-rate-district-scatterplot}Relationship
between the sum of contributions over \$200 donated in- v.s.
out-of-district, for each district during the 2013 representational
order (effective between 2015 and 2024). (A) shows this relationship for
donations to EDAs, and (B) shows this for donations to candidates.
Points represent individual districts. The grey dashed line is 1:1. The
top 3 districts by amount contributed in- and out-of-district
respectively are labelled with their names. Graphs are similar when
counts are displayed instead of amounts. (A) is similar when only
donations sent during writ periods are considered, and (B) is similar
when specific election cycles are considered.}

\end{figure}%

\subsubsection{Characterizing Inter-District Funding
Flows}\label{sec-data-summaries-ood-network}

Finally, we briefly describe some characteristics of the localized
funding network formed by out-of-district donations.
Figure~\ref{fig-data-summaries-ood-network-senders} displays the
relationship between the amount donated out-of-district and the amount
received from out-of-district, for each district. For both entities, the
majority of districts fall cleanly into two groups; those that send more
than they receive (e.g.~Trinity --- St.~Paul's), and those that receive
more than they send (e.g.~Scarborough Centre).

Table~\ref{tbl-data-summaries-ood-neighbours} shows the proportion of
contributions that are received by neighbouring districts, and by
districts out-of-province, among all out-of-district contributions.
These proportions offer a rough picture of how far the average
out-of-district donation travels\footnote{A distribution of donation
  distances (measured from the centroid of donor's dissemination area to
  the centroid of the recipient's district) is displayed in
  Appendix~\ref{sec-district-propensity-simulation}.}. For both
entities, around 40\% of out-of-district donations are sent to adjacent
districts, while only around 12\% reach other provinces. Out-of-district
donations are overall closer to home for candidate-bound donations,
echoing our findings in Table~\ref{tbl-data-summaries-ood-rate-overall}.
Those out-of-province donations are visualized further in
Figure~\ref{fig-data-summaries-ood-network-provinces}. For donations to
district associations, the total flux (sum of contributed amount
received and sent) of a province strongly correlates with its
population, and contribution rates are larger between two highly
populated provinces. The network of interprovincial funding of
candidates, however, has different primary edges, and overall
contributions are more concentrated between B.C., Alberta, and Ontario.
While we should expect reduced interprovincial contribution rates to
candidates compared to EDAs
(Table~\ref{tbl-data-summaries-ood-neighbours}), certain differences in
the candidate network such as Ontario and B.C. concentrating their
outflow to Alberta, represent genuine differences in interprovincial
funding regimes between these entities.

\begin{figure}

\centering{

\pandocbounded{\includegraphics[keepaspectratio]{paper_files/figure-pdf/fig-data-summaries-ood-network-senders-1.pdf}}

}

\caption{\label{fig-data-summaries-ood-network-senders}Relationship
between the sum of contributions over \$200 sent out-of-district v.s.
received from out-of-district, for each district during the 2013
representational order (effective between 2015 and 2024). (A) shows this
relationship for donations to EDAs, and (B) shows this for donations to
candidates. Points represent individual districts. The grey dashed line
is 1:1. The top three districts by amount sent and received respectively
are labelled with their names. Graphs are similar when counts are
displayed instead of amounts. (A) is similar when only donations sent
during writ periods are considered, and (B) is similar when specific
election cycles are considered.}

\end{figure}%

\begin{table}

\caption{\label{tbl-data-summaries-ood-neighbours}Proportion of
out-of-district contributions over \$200 that were sent to adjacent
districts, and out-of-province during our study period. Raw and
value-weighted proportions are reported, along with the overall number
of contributions for reference. Contributions received by candidates
outside general election cycles (by-elections) are not included in
calculations.}

\centering{

\centering
\begin{tblr}[         %% tabularray outer open
]                     %% tabularray outer close
{                     %% tabularray inner open
width={1\linewidth},
colspec={X[]X[]X[]X[]},
hline{2}={1-4}{solid, black, 0.05em},
hline{1}={1-4}{solid, black, 0.08em},
hline{6}={1-4}{solid, black, 0.08em},
cell{1}{1}={}{font=\bfseries\fontsize{0.85em}{1.15em}\selectfont, halign=l},
cell{1}{2}={}{font=\bfseries\fontsize{0.85em}{1.15em}\selectfont, halign=l},
cell{1}{3}={}{font=\bfseries\fontsize{0.85em}{1.15em}\selectfont, halign=r},
cell{1}{4}={}{font=\bfseries\fontsize{0.85em}{1.15em}\selectfont, halign=r},
cell{2-5}{1}={}{font=\fontsize{0.85em}{1.15em}\selectfont, halign=l},
cell{2-5}{2}={}{font=\fontsize{0.85em}{1.15em}\selectfont, halign=l},
cell{2-5}{3}={}{font=\fontsize{0.85em}{1.15em}\selectfont, halign=r},
cell{2-5}{4}={}{font=\fontsize{0.85em}{1.15em}\selectfont, halign=r},
}                     %% tabularray inner close
Received by & Sent to & By count & By amount \\
Candidate & Adjacent district & 43.3\% & 44.1\% \\
Candidate & Out-of-province & 10.6\% & 9.7\% \\
District association & Adjacent district & 37.1\% & 36.0\% \\
District association & Out-of-province & 14.7\% & 14.4\% \\
\end{tblr}

}

\end{table}%

\begin{figure}

\centering{

\pandocbounded{\includegraphics[keepaspectratio]{paper_files/figure-pdf/fig-data-summaries-ood-network-provinces-1.pdf}}

}

\caption{\label{fig-data-summaries-ood-network-provinces}Interprovincial
funding flows for contributions over \$200 during the 2013
representational order (effective between 2015 and 2024). (A) shows
flows for donations to EDAs, and (B) shows flows for donations to
candidates. Networks were similar when counts were displayed instead of
amounts. Node size reflects total interprovincial flux (amount sent +
amount received), and edges are weighted by transfer amount. Provinces
with red-coloured nodes receive more contributions than they send, and
vice-versa for blue nodes. Edges are coloured on a gradient from blue
(sending province) to red (receiving province) to mark direction.
Candidate networks for individual writ-periods were similar to their
aggregate displayed here.}

\end{figure}%

\subsection{Out-of-District-Propensity
Models}\label{sec-model-specification}

\subsubsection{Individual-Level
Model}\label{sec-model-specification-individual}

We used a logistic regression model to explore how the likelihood of
out-of-district donation varies across different transactional
characteristics, donor demographic environments, and district
competitiveness, for individual donations. Our outcome variable is
binary: whether or not a donation was sent out-of-district. This model
was fit on both donations to EDAs and candidates combined. Models fit on
each localized entity separately had similar coefficient magnitudes,
signs and significances (see
Appendix~\ref{sec-alternate-model-specifications}). Donations were not
grouped by individual, since we were interested in both donor and
recipient characteristics in this model. Diagnostics showed that
residual autocorrelation, the primary statistical risk associated with
this decision, was not a large concern in the model.

The transactional predictors in our model are: the type of recipient
(candidate or district association), their political party, and the
monetary value of the donation. We did not include time-indexed
variables since our data summaries showed that out-of-district rates
were constant over time, for both candidates and district associations.
Following Garnett (\citeproc{ref-garnett_where_2024}{2024}), we also
considered the following characteristics of the donor's
socio-demographic environment: province (grouped into regions), mean
age, median household income, the share of the population with a
post-secondary degree, and the share of the population that is a visible
minority\footnote{An alternative model where regions were disaggregated
  into provinces showed qualitatively similar results. An alternative
  model where the visible minority proportion was disaggregated into
  specific ethnic groups is shown in
  Appendix~\ref{sec-alternate-model-specifications}.}. These predictors
are associated with out-of-district donation in the U.S.
(\citeproc{ref-gimpel_political_2006}{Gimpel et al. 2006},
\citeproc{ref-gimpel_check_2008}{2008}). These environmental
characteristics were assigned to donations based on their donor's
dissemination area. District competitiveness was proxied using the
difference between the percentage of votes received by the winning and
runner-up candidates, averaged across all three general elections
captured in our study period\footnote{Districts were more likely to have
  elected an MP from the same party (or simply the same MP) for all
  three elections when their average margin of victory is higher
  (\(\text{pseudo} \space R^2 = 0.22\)), meaning that this
  competitiveness metric also captures a decent share of the variation
  in district-level party stability. These data were sourced from Open
  Government Canada (\citeproc{ref-canada_open_government_2026}{Treasury
  Board of Canada Secretariat 2026}).}, inspired by Garnett
(\citeproc{ref-garnett_where_2024}{2024}). This predictor was included
since we might expect donors from less competitive districts to donate
out-of-district more often, recognizing that a local contribution will
do little to affect an outcome that is effectively predetermined.

We also included the population density of the donor's district, for two
reasons. First, out-of-district donors are more likely to come from
urban (and therefore more densely populated) districts in the U.S.
(\citeproc{ref-gimpel_check_2008}{Gimpel et al. 2008}). Second,
population density serves as a relatively good proxy for how spatially
clustered a district is, since Canadian districts are designed to have
similar population sizes
(\citeproc{ref-canada_redistribution_2025}{Elections Canada 2025b}).
Out-of-district donation rates are likely naturally higher in smaller
districts, since donors are more likely to interact with individuals,
businesses, or politicians close by that happen to be located in a
different district. Population density was transformed in the model by
applying a square root, allowing us to interpret the metric as the
inverse of the average distance between people, which captures this
distance effect more faithfully\footnote{The observed spatial variation
  in out-of-district propensity cannot be fully explained by the spatial
  arrangement of districts alone
  (Appendix~\ref{sec-district-propensity-simulation}), however it
  remains an important consideration.}.

The resulting model is:

\[
\begin{aligned}
\text{logit}(p_i) &= \beta_0 + \beta_1 \text{Election}_i + \beta_2 \text{Candidate}_{i} + \sum_{j=3}^{5} \beta_j \text{Party}_{ji} + \beta_6 \text{Amount}_i \\
&\quad + \sum_{k=7}^{11} \beta_k \text{Region}_{ki} + \beta_{12} \text{Age}_i + \beta_{13} \text{Income}_i + \beta_{14} \text{Education}_i + \beta_{15} \text{VisMin}_i \\
&\quad + \beta_{16} \text{VictoryMargin}_i + \beta_{17} \text{Density}_i \\
\end{aligned}
\]

where:

\begin{itemize}
\tightlist
\item
  \(p_i\) is the probability that donation \(i\) was an out-of-district
  donation.
\item
  \(\text{Election}_i\) indicates whether donation \(i\) occurred during
  a general election period.
\item
  \(\text{Candidate}_{i}\) indicates whether donation \(i\) was
  addressed to a Candidate (baseline is EDA).
\item
  \(\text{Party}_{ji}\) indicates which party (among Liberal, N.D.P, or
  Other, with baseline being Conservative) the recipient is associated
  with.
\item
  \(\text{Amount}_i\) is the monetary value of donation \(i\).
\item
  \(\text{Region}_{ki}\) indicates which region (among Atlantic Canada,
  B.C., Prairies, Quebec, Territories, with baseline being Ontario) the
  donor is located.
\item
  \(\text{Age}_i\) is the average age of residents in the donor's
  dissemination area.
\item
  \(\text{Income}_i\) is the median household income in the donor's
  dissemination area.
\item
  \(\text{Education}_i\) is the proportion of residents in the donor's
  home district holding a post-secondary qualification.
\item
  \(\text{VisMin}_i\) is the proportion of residents in the donor's home
  district who identify as a visible minority.
\item
  \(\text{VictoryMargin}_i\) is the difference between the percentage of
  votes received by the winning and runner-up candidates, averaged
  across all three general elections in the donor's home district.
\item
  \(\text{Density}_i\) is the population density (persons per \(km^2\))
  of the donor's home district. As mentioned in-text, the model was fit
  using the square root of this variable.
\end{itemize}

\subsubsection{District-Level
Model}\label{sec-model-specification-district}

Inspired by the large variation in out-of-district giving between
districts observed in
Figure~\ref{fig-data-summaries-ood-rate-district-scatterplot}, we fit a
second regression model to explore which socio-demographic features were
associated with higher rates of out-of-district donation, aggregated at
the district level. Our outcome variable in this case was the proportion
of the total contributed amount\footnote{We observed high correlation
  (\(\rho = 0.978\)) between the proportion of out-of-district donations
  and the share of total contributed amount that left the district, for
  both candidates and EDAs. Therefore, results would have been similar
  if we had used the number of donations as our outcome instead of the
  amount-weighted share of donations.} sent out-of-district during our
entire study period, for each district. Given that this outcome is a
proportion, we used beta regression, fit using the \texttt{betareg}
package in R (\citeproc{ref-betareg}{Cribari-Neto and Zeileis 2010}).
This model was fit on both donations to EDAs and candidates combined,
following from the individual model. A model fit on only EDAs separately
had similar coefficient magnitudes, and the same signs and significances
(shown in Appendix --~\ref{sec-alternate-model-specifications}). A model
fit on only candidates was not fit due to data sparsity and the presence
of singularities in the response.

The covariates included in this model were the same environment
variables chosen for the individual-level model, aggregated at the
district level\footnote{As with the individual-level model, alternative
  model specifications with where disaggregated region and visible
  minority proportion is shown in
  Appendix~\ref{sec-alternate-model-specifications}.}. We also included
the same district-level competitiveness metric used in the individual
model. Since districts had very different numbers of donations used to
calculate our outcome proportion, we also assumed a power-law
relationship between the common precision parameter \(\phi\) and the
number of donations, following Simas et al.
(\citeproc{ref-generalized_beta_regression}{2010}).

The resulting model is:

\[
\begin{aligned}
p_i &\sim \text{Beta}(\mu_i, \phi_i) \\
\text{logit}(\mu_i) &= \beta_0 + \beta_1 \text{Amount}_i + \sum_{j=2}^{6} \beta_j \text{Region}_{ji} + \beta_7 \text{Age}_i + \beta_8 \text{Income}_i \\
&\quad  + \beta_9 \text{Education}_i + \beta_{10} \text{VisMin}_i + \beta_{11} \text{VictoryMargin}_i + \beta_{12} \text{Density}_i \\
\log(\phi_i) &= \gamma_0 + \gamma_1\log(n_i)
\end{aligned}
\]

where:

\begin{itemize}
\tightlist
\item
  \(p_i\) is the proportion of the total contributed amount that was
  sent out-of-district by donors in district \(i\).
\item
  \(\mu_i\) and \(\phi_i\) are the mean and common precision parameters
  of the distribution of \(p_i\).
\item
  \(n_i\) is the number of donations originating from district \(i\).
\item
  \(\text{Amount}_i\) is the total amount contributed by donors in
  district \(i\).
\item
  \(\text{Region}_{ji}\) are indicators for the donor region of district
  \(i\).
\item
  \(\text{Age}_i\) is the mean age of residents in district \(i\).
\item
  \(\text{Income}_i\) is the median household income in district \(i\).
\item
  \(\text{Education}_i\) is the proportion of residents in district
  \(i\) holding a post-secondary qualification.
\item
  \(\text{VisMin}_i\) is the proportion of residents in district \(i\)
  who identify as a visible minority.
\item
  \(\text{VictoryMargin}_i\) is the difference between the percentage of
  votes received by the winning and runner-up candidates, averaged
  across all three general elections in district \(i\).
\item
  \(\text{Density}_i\) is the population density (persons per \(km^2\))
  of district \(i\). As mentioned in-text, the model was fit using the
  square root of this variable.
\end{itemize}

\section{Results}\label{sec-results}

Here we present the results from our out-of-district propensity models.
We note that many of our key findings are best understood as data
summaries and visualizations, which are presented in
Section~\ref{sec-data-summaries}.
Section~\ref{sec-data-summaries-situation} provides context for how the
localized contributions considered in our study fit into the broader
scope of financial contributions in general.
Section~\ref{sec-data-summaries-ood-rate} provides spatio-temporal
summaries of out-of-district propensity (which were used to inform the
model discussed here). Section~\ref{sec-data-summaries-ood-network}
briefly characterizes the funding network formed by out-of-district
donations. Both summary-based and model-based findings are discussed
together in Chapter~\ref{sec-discussion}.

Regression results from the individual- and district-level models are
presented in Table~\ref{tbl-results-model-individual} and
Table~\ref{tbl-results-model-district} respectively.
Figure~\ref{fig-results-models-coefs-individual} and
Figure~\ref{fig-results-models-coefs-district} provide a visualization
of coefficient signs, magnitudes, and significances. For both models,
exponentiating the coefficients yields a more interpretable value, the
odds ratio. For the individual (i.e.~logistic) regression model, an odds
ratio above one for a categorical variable (e.g.~recipient belonging to
the Liberal Party) indicates that, holding all else constant, donations
of that category have increased odds of being out-of-district, relative
to the baseline category (e.g.~recipient belonging to the Conservative
party). For the district level (i.e.~beta) regression model, an odds
ratio above one for a categorical variable (e.g.~district situated in
the Prairies) indicates that, holding all else constant, the expected
ratio of out- vs.~in-district contributions increases. Odds ratios less
than one and continuous variables (e.g.~median household income) are
interpreted accordingly.

Individual donations are less likely to reach out-of-district when they
are bound for candidates instead of district associations (EDAs), and
more likely for Liberal recipients instead of Conservative recipients.
The contribution likelihoods of Conservative and N.D.P recipients were
not meaningfully different. Larger donations were more likely to leave
the district compared to smaller ones, helping to explain why
amount-weighted proportions in our data summaries were consistently
larger. The average margin of victory (taken to be a measure of district
non-competitiveness) was, expectedly, positively associated with
out-of-district propensity. Donations originating from all regions were
less likely to reach out-of-district compared to Ontario, with the
exception of Quebec, whose donors have nearly 1.5 times the odds to
donate out-of-district, holding all else constant. Donations originating
in the Territories were especially unlikely to travel outside their
district, with nearly one-eighth the baseline Ontario odds. Turning to
the donor's socio-demographic environment, regions with higher
population densities, post-secondary education rates, median household
incomes, and visible minority proportions\footnote{This effect was
  observed even when donations from the four Brampton constituencies
  were removed, since data summaries showed that these clearly had very
  high influence on the model fit. A similar check was performed on the
  district-level model.} all favoured out-of-district donation. However,
the average age did not significantly impact out-of-district chances in
the presence of other socio-demographic variables\footnote{Average age
  was significant in individual and district-level models where it,
  along with population density, were the only socio-demographic
  predictors. This tells us that its effect is mostly explained by other
  factors, such as income or educational attainment.}. When ethnicities
were disaggregated, the positive effect of the visible minority
proportion was observed for all minority ethnic groups recorded in the
Census (see Appendix~\ref{sec-alternate-model-specifications}).

At the district level, we modelled the ratio of out- to in-district
contribution sums. Similarly to the individual models, districts with
higher population densities, post-secondary education rates, median
household incomes, and visible minority proportions were more likely to
allocate more of their contributions to political entities in other
districts rather than their own. The average age of a district was,
similarly, insignificant in the presence of other socio-demographic
variables. Quebec and the Territories were the only regions to differ
significantly from Ontario in their out-of-district ratios, and with the
same coefficient signs as in the individual model. Unlike the individual
model, the total amount of donations raised and the average margin of
victory in a district did not have a significant effect on the
proportion of that amount that left said district; while donations are
more likely to travel out-of-district when they are larger, and when
districts are less competitive, these effects seems to be a purely
individualized.

The eight variables used in our district-level model were able to
explain a large amount of the variance in the ratio of out- to
in-district donation (\({pseudo} \space R^2 = 0.67\)). The same model
fit on only population density also achieved a relatively good fit
(\({pseudo} \space R^2 = 0.45\)), however other variables, and
especially socio-demographic ones such as educational attainment and
visible minority proportion, improved the model's explanatory power well
past what population density alone could offer.

\begin{table}

\caption{\label{tbl-results-model-individual}Results from a logistic
regression model where the outcome is whether or not an individual
donation was sent out-of-district. The reference entity is the district
association, the reference party is the Conservative Party, and the
reference region is Ontario. Estimates are bolded when significant at
the alpha = 0.05 level. Small losses (\textasciitilde{} 1.1\% of
entries) were incurred due to missingness in the census data.}

\centering{

\centering
\begin{tblr}[         %% tabularray outer open
]                     %% tabularray outer close
{                     %% tabularray inner open
colspec={Q[]Q[]Q[]Q[]Q[]},
hline{2}={5}{solid, black, 0.03em},
hline{2}={2,4}{solid, black, 0.03em, l=-0.5},
hline{2}={3}{solid, black, 0.03em, r=-0.5},
hline{3}={1-5}{solid, black, 0.05em},
hline{23}={1-5}{solid, black, 0.05em},
hline{1}={1-5}{solid, black, 0.08em},
hline{29}={1-5}{solid, black, 0.08em},
cell{1}{1}={}{halign=l},
cell{1}{2}={c=2}{halign=c},
cell{1}{3}={}{halign=c},
cell{1}{4}={c=2}{halign=c},
cell{1}{5}={}{halign=c},
cell{2-3,5,7,9-10,12-16,18-22}{2}={}{font=\bfseries\fontsize{0.85em}{1.15em}\selectfont, halign=c},
cell{2-3,5,7,9-10,12-16,18-22}{3}={}{font=\bfseries\fontsize{0.85em}{1.15em}\selectfont, halign=c},
cell{2-3,5,7,9-10,12-16,18-22}{4}={}{font=\bfseries\fontsize{0.85em}{1.15em}\selectfont, halign=c},
cell{2-3,5,7,9-10,12-16,18-22}{5}={}{font=\bfseries\fontsize{0.85em}{1.15em}\selectfont, halign=c},
cell{2}{1}={}{font=\bfseries\fontsize{0.85em}{1.15em}\selectfont, halign=l},
cell{3-28}{1}={}{font=\fontsize{0.85em}{1.15em}\selectfont, halign=l},
cell{4,6,8,11,17,23-28}{2}={}{font=\fontsize{0.85em}{1.15em}\selectfont, halign=c},
cell{4,6,8,11,17,23-28}{3}={}{font=\fontsize{0.85em}{1.15em}\selectfont, halign=c},
cell{4,6,8,11,17,23-28}{4}={}{font=\fontsize{0.85em}{1.15em}\selectfont, halign=c},
cell{4,6,8,11,17,23-28}{5}={}{font=\fontsize{0.85em}{1.15em}\selectfont, halign=c},
}                     %% tabularray inner close
& Coefficient &  & Odds Ratio &  \\
& Est. & S.E. & Est. & S.E. \\
(Intercept) & \num{-2.889} & \num{0.062} & \num{0.056} & \num{0.003} \\
District Associations & - &  & - &  \\
Candidates & \num{-0.150} & \num{0.015} & \num{0.861} & \num{0.013} \\
Conservative & - &  & - &  \\
Liberal & \num{0.319} & \num{0.014} & \num{1.376} & \num{0.020} \\
N.D.P & \num{0.013} & \num{0.018} & \num{1.013} & \num{0.019} \\
Other & \num{0.349} & \num{0.028} & \num{1.418} & \num{0.040} \\
Amount & \num{0.000} & \num{0.000} & \num{1.000} & \num{0.000} \\
Ontario & - &  & - &  \\
Quebec & \num{0.318} & \num{0.020} & \num{1.375} & \num{0.028} \\
British Columbia & \num{-0.431} & \num{0.018} & \num{0.650} & \num{0.011} \\
Prairies & \num{-0.386} & \num{0.019} & \num{0.680} & \num{0.013} \\
Atlantic Canada & \num{-0.262} & \num{0.027} & \num{0.769} & \num{0.021} \\
Territories & \num{-1.995} & \num{0.145} & \num{0.136} & \num{0.020} \\
Age & \num{-0.002} & \num{0.001} & \num{0.998} & \num{0.001} \\
Income & \num{0.000} & \num{0.000} & \num{1.000} & \num{0.000} \\
Education & \num{1.923} & \num{0.059} & \num{6.841} & \num{0.406} \\
Visible Minority & \num{2.205} & \num{0.029} & \num{9.071} & \num{0.260} \\
Margin of Victory (mean) & \num{0.007} & \num{0.001} & \num{1.007} & \num{0.001} \\
Population Density (sqrt) & \num{0.003} & \num{0.000} & \num{1.003} & \num{0.000} \\
Num.Obs. & \num{136437} &  & \num{136437} &  \\
AIC & \num{164482.3} &  & \num{164482.3} &  \\
BIC & \num{164649.3} &  & \num{164649.3} &  \\
Log.Lik. & \num{-82224.138} &  & \num{-82224.138} &  \\
F & \num{1115.752} &  & \num{1115.752} &  \\
RMSE & \num{0.46} &  & \num{0.46} &  \\
\end{tblr}

}

\end{table}%

\begin{table}

\caption{\label{tbl-results-model-district}Results from a beta
regression model where the outcome is the out-of-district share of the
sum of contributions over \$200 originating from an individual district
during our study period. The reference entity is the district
association, and the reference region is Ontario. Estimates are bolded
when significant at the alpha = 0.05 level. Small losses
(\textasciitilde{} 1.1\% of individual transactions) were incurred due
to missingness in the census data. Each of the 338 districts was
represented in this model.}

\centering{

\centering
\begin{tblr}[         %% tabularray outer open
]                     %% tabularray outer close
{                     %% tabularray inner open
colspec={Q[]Q[]Q[]Q[]Q[]},
hline{2}={5}{solid, black, 0.03em},
hline{2}={2,4}{solid, black, 0.03em, l=-0.5},
hline{2}={3}{solid, black, 0.03em, r=-0.5},
hline{3}={1-5}{solid, black, 0.05em},
hline{17}={1-5}{solid, black, 0.05em},
hline{1}={1-5}{solid, black, 0.08em},
hline{22}={1-5}{solid, black, 0.08em},
cell{1}{1}={}{halign=l},
cell{1}{2}={c=2}{halign=c},
cell{1}{3}={}{halign=c},
cell{1}{4}={c=2}{halign=c},
cell{1}{5}={}{halign=c},
cell{2-3,9-10,12-14,16}{2}={}{font=\bfseries\fontsize{0.85em}{1.15em}\selectfont, halign=c},
cell{2-3,9-10,12-14,16}{3}={}{font=\bfseries\fontsize{0.85em}{1.15em}\selectfont, halign=c},
cell{2-3,9-10,12-14,16}{4}={}{font=\bfseries\fontsize{0.85em}{1.15em}\selectfont, halign=c},
cell{2-3,9-10,12-14,16}{5}={}{font=\bfseries\fontsize{0.85em}{1.15em}\selectfont, halign=c},
cell{2}{1}={}{font=\bfseries\fontsize{0.85em}{1.15em}\selectfont, halign=l},
cell{3-21}{1}={}{font=\fontsize{0.85em}{1.15em}\selectfont, halign=l},
cell{4-8,11,15,17-21}{2}={}{font=\fontsize{0.85em}{1.15em}\selectfont, halign=c},
cell{4-8,11,15,17-21}{3}={}{font=\fontsize{0.85em}{1.15em}\selectfont, halign=c},
cell{4-8,11,15,17-21}{4}={}{font=\fontsize{0.85em}{1.15em}\selectfont, halign=c},
cell{4-8,11,15,17-21}{5}={}{font=\fontsize{0.85em}{1.15em}\selectfont, halign=c},
}                     %% tabularray inner close
& Coefficient &  & Odds Ratio &  \\
& Est. & S.E. & Est. & S.E. \\
(Intercept) & \num{-4.276} & \num{0.797} & \num{0.014} & \num{0.011} \\
Amount & \num{-0.000} & \num{0.000} & \num{1.000} & \num{0.000} \\
Ontario & - &  & - &  \\
Quebec & \num{0.556} & \num{0.112} & \num{1.744} & \num{0.196} \\
British Columbia & \num{-0.166} & \num{0.098} & \num{0.847} & \num{0.083} \\
Prairies & \num{-0.144} & \num{0.109} & \num{0.866} & \num{0.094} \\
Atlantic Canada & \num{0.228} & \num{0.138} & \num{1.256} & \num{0.173} \\
Territories & \num{-1.339} & \num{0.483} & \num{0.262} & \num{0.126} \\
Age & \num{0.015} & \num{0.016} & \num{1.015} & \num{0.016} \\
Income & \num{0.000} & \num{0.000} & \num{1.000} & \num{0.000} \\
Education & \num{1.654} & \num{0.600} & \num{5.226} & \num{3.137} \\
Visible Minority & \num{1.918} & \num{0.221} & \num{6.805} & \num{1.505} \\
Margin of Victory (mean) & \num{0.004} & \num{0.003} & \num{1.004} & \num{0.003} \\
Population Density (sqrt) & \num{0.012} & \num{0.002} & \num{1.012} & \num{0.002} \\
Num.Obs. & \num{338} &  & \num{338} &  \\
R2 & \num{0.672} &  & \num{0.672} &  \\
AIC & \num{-473.8} &  & \num{-473.8} &  \\
BIC & \num{-416.4} &  & \num{-416.4} &  \\
RMSE & \num{0.13} &  & \num{0.13} &  \\
\end{tblr}

}

\end{table}%

\begin{figure}

\centering{

\pandocbounded{\includegraphics[keepaspectratio]{paper_files/figure-pdf/fig-results-models-coefs-individual-1.pdf}}

}

\caption{\label{fig-results-models-coefs-individual}Coefficient
estimates and 95\% confidence intervals from a logistic regression model
where the outcome is whether or not an individual donation was sent
out-of-district. Coefficients are coloured grey if they are
insignificant at alpha = 0.05, red if they are significant with a
negative effect, and green if they are significant with a positive
effect. Reference entity is the district association, reference party is
the Conservative Party, reference region is Ontario.}

\end{figure}%

\begin{figure}

\centering{

\pandocbounded{\includegraphics[keepaspectratio]{paper_files/figure-pdf/fig-results-models-coefs-district-1.pdf}}

}

\caption{\label{fig-results-models-coefs-district}Coefficient estimates
and 95\% confidence intervals from a beta regression model where the
outcome is the out-of-district share of the sum of contributions over
\$200 originating from an individual district during our study period.
Coefficients are coloured grey if they are insignificant at alpha =
0.05, red if they are significant with a negative effect, and green if
they are significant with a positive effect. Reference entity is the
district association, reference region is Ontario.}

\end{figure}%

\section{Discussion}\label{sec-discussion}

In this paper, we used publicly available political financing data to
document the extent and spatio-temporal trends in Canadian
out-of-district donations from 2015 to 2024. Further, we identified
donor-environmental characteristics associated with out-of-district
propensity (i.e.: where out-of-district donations are more or less
likely), and provided a preliminary overview of the spatial network
formed by out-of-district donations. Our findings are also
contextualized within the broader scope of all\footnote{Excluding a
  small number of contributions (1.2\%) where we were not able to
  geolocate the donor.} Canadian federal political contributions over
\$200 during this time (see Section~\ref{sec-data-summaries-situation}).
We present and discuss our major findings here, along with some notable
study limitations.

Out-of-district contributions constituted a substantial source of
funding during our study period, to electoral campaigns and district
associations alike. Naturally, overall contribution rates to these two
political entities varied in time; campaign funding is largely
restricted to writ periods, and donations to district associations
swelled during the holiday season and electoral cycles, which garnered
nearly five times more contributions over \$200 compared to baseline
rates. However, for both entities, the pooled share of funding received
from outside their district (equivalently, the share of the total
contributed amount that left its donor's district) remained largely
constant throughout the decade. Roughly 40\% of individual donations
over \$200 (accounting for 44\% of the total amount contributed) are
sent out-of-district, a rate which remained stable throughout our study
period. This rate and its consistency through time contrasts against
U.S. findings, where the bulk of existing out-of-district donations
research has been done. There, out-of-district funding shares to
senators and congressional representatives during electoral campaigns
have increased since the start of the 21st century
(\citeproc{ref-gimpel_political_2006}{Gimpel et al. 2006};
\citeproc{ref-canes_wrone_out_2022}{Canes-Wrone and Miller 2022};
\citeproc{ref-sievert_out_2023}{Sievert and Mathiasen 2023}), and the
average representative in 2016 received around two-thirds of their
contributions from out-of-district
(\citeproc{ref-canes_wrone_out_2022}{Canes-Wrone and Miller 2022}).

While out-of-district contribution rates varied little over time, across
parties, and between the recipient entity types that we considered
(candidates or district associations), we document drastic differences
in out-of-district propensity from district-to-district over our study
period, both at the level of individual donations, and entire districts.
For example, out-of-district donation shares to district associations
over the entire studied decade were as little as 5\% for some districts
(e.g.~Niagara Falls), and as high as 75\% in others (e.g.~Trinity ---
St.~Paul's). While we might expect passively higher out-of-district
rates solely due to the spatial proximity of other districts, empirical
variation in rates of out-of-district giving (aggregated at the district
level) far exceeded simulated expectations under this hypothesis.
Further, most deviations from expected out-of-district donation rates
were districts with higher, not lower observed out-of-district
contribution shares (see
Appendix~\ref{sec-district-propensity-simulation}). This variation is
also not simply a result of the large differences in the number of
donations receive by different districts; restricting our sample to only
those districts among the top 10\% by total contributed amount yielded a
similarly broad range of out-of-district propensities.

Donations originating from densely populated, high-income, and
well-educated districts are more likely to travel out-of-district, all
else held constant. This set of covariates agrees with findings
established in the U.S. (\citeproc{ref-gimpel_check_2008}{Gimpel et al.
2008}). These are also significant predictors of donation rates in
general, across North America (\citeproc{ref-gimpel_check_2008}{Gimpel
et al. 2008}; \citeproc{ref-garnett_where_2024}{Garnett 2024}). In
addition to these, however, the visible minority proportion of a donor's
home environment is positively associated with out-of-district giving,
and was the single most explanatory socio-demographic predictor in both
of our models. This finding departs from the established profile of
American ``donor districts'', which historically have predominantly
European ethnic makeups (\citeproc{ref-gimpel_check_2008}{Gimpel et al.
2008}). However, recent Canadian studies support this inverted
relationship between ethnicity and political engagement. Garnett et al.
(\citeproc{ref-garnett_lifeblood_2022}{2022}), among others, note that
some minority populations such as South Asians donate at rates that far
exceed their share of the population, and, topically, Besco and Tolley
(\citeproc{ref-besco_ethnic_2022}{2022}) found that individual donors
belonging to most minority ethnic groups show higher rates of
out-of-district donation compared to donors of European descent. The
strong relationship between highly minoritized regions of the country
and rates of out-of-district giving identified in our study is most
plausibly attributed to one of two mechanisms which our models, as
specified, cannot disentangle; demographic collinearity, and individual
co-ethnic affinity. First, the relationship may simply be a byproduct of
Canada's unique social demography. Visible minority populations in
Canada are concentrated almost exclusively in cities, and higher
minority composition is not obviously associated with lower income (a
well-established predictor of donation and out-of-district donation
propensity) at the national level (\citeproc{ref-hou_recent_2004}{Hou
2004}; \citeproc{ref-walks_ghettos_2006}{Walks and Bourne 2006}). In the
U.S., however, areas with high visible minority proportions do not
follow an urban-rural gradient nearly as closely
(\citeproc{ref-lee_ethnoracial_2017}{Lee and Sharp 2017}) and do have
lower income, on average
(\citeproc{ref-quillian_segregation_2012}{Quillian 2012}). Therefore, it
is entirely possible that our visible minority coefficient is absorbing
residual features of affluent, dense, urban Canada that density and
income alone do not fully capture, rather than reflecting anything about
ethnicity itself. More generally, collinearity among socio-demographic
predictors was present in our dataset, making it difficult to
disentangle the specific drivers of out-of-district donation, even
though it is undeniable that social demography and out-of-district
donation are linked, beyond simple population density. Alternatively,
this relationship may be signalling the presence of co-ethnic
inter-district donor networks. Identity-based donor affinity is a
well-established feature of Canadian political participation
(\citeproc{ref-besco_identities_2019}{Besco 2019};
\citeproc{ref-besco_ethnic_2022}{Besco and Tolley 2022};
\citeproc{ref-tolley_who_2022}{Tolley et al. 2022}). Further, the
largest component of our inter-district donation network we could
identify is the flow of funds between the four districts that make up
the Brampton region within the Greater Toronto Area, whose population is
majority South Asian. While attractive, this hypothesis cannot be
confirmed without individualized ethnic identification data, largely
inaccessible at this scale.

Despite empirically high rates of out-of-district donation, our study
provides evidence that these donations remain a largely regionalized
phenomenon. Nearly half of donations that leave their district during
our study period were destined for entities in adjacent ones, and only
around 12\% ever even left their donor's province (specifically, 9\% for
candidate-bound donations). U.S. research on out-of-state donations
describes highly nationalized donor bases among incumbent and
campaigning House representatives
(\citeproc{ref-gimpel_political_2006}{Gimpel et al. 2006};
\citeproc{ref-canes_wrone_out_2022}{Canes-Wrone and Miller 2022};
\citeproc{ref-sievert_out_2023}{Sievert and Mathiasen 2023};
\citeproc{ref-keena_out_2024}{Keena 2024}). For example, a study on
senatorial elections from 2000 to 2020 found that both Republican and
Democrat senators received at least 49\% of their donations from
out-of-state during election seasons
(\citeproc{ref-sievert_out_2023}{Sievert and Mathiasen 2023}). Some of
this difference may simply be a matter of spatial proximity; U.S. states
are closer to each other than Canadian provinces and territories. The
scale of this difference, however, speaks to a more fundamental
difference between the geographical donor bases of each country;
Canadian electoral candidates largely source individual-based funding
from within their own province. For those few donations bound for other
provinces, we observe qualitatively different patterns for
candidate-bound and district association-bound donations. Further
research on Canadian out-of-district donations may take advantage of the
spatial network they form to explore these patterns further.

Our findings are accompanied by some important limitations. First, a
non-negligible proportion (34\%) of the full amount addressed to
localized entities were under \$200, and therefore not included in our
study. The majority (75\%) of these donations were location-anonymized,
and were excluded out of necessity, following Gimpel et al.
(\citeproc{ref-gimpel_check_2008}{2008}) and Garnett
(\citeproc{ref-garnett_where_2024}{2024}). Distributional comparisons
between study and location-anonymized contributions (see
Appendix~\ref{sec-missing-data}) showed that this missing sum was
distributed similarly to the study data across recipient characteristics
such as party affiliation, region, and over time, suggesting that their
exclusion introduced little systematic bias into our study dataset.
However, location-disclosed contributions under \$200 were distributed
very differently to those above \$200 included in our study: to an
overwhelming degree, these were in-district donations to Liberal
district associations outside of election cycles, and were considerably
less likely to be sent out-of-district than those in our study (18\%
compared to 42\%). Further investigation into these cases revealed that
only 12.4\% of these transactions had a unique donor name, postal code,
and year, compared to 92.7\% uniqueness across the same set of factors
in our study data. This tells us that a majority of these transactions
are likely from donors who, in a year, contributed a total above \$200
to district associations, and therefore had their locations disclosed.
Given this, the exclusion of location-disclosed contributions under
\$200 likely inflates our estimates of out-of-district rates for Liberal
district associations, which were already found to receive the largest
contribution share from out-of-district among all entity-party pairs
(57\% compared to the average of 40\%). More broadly, this finding,
paired with the positive association between amount donated and
out-of-district propensity, suggests that our estimates of
out-of-district donation rates, particularly the amount-unweighted point
estimates, are overestimates of the true underlying out-of-district
propensity. Nevertheless, since contributions under \$200 remained
similarly distributed to the study data, particularly in space, their
exclusion is unlikely to have greatly affected the statistical
significance of our most interesting covariates in our out-of-district
propensity models, including those related to social demography.

Second, our focus on out-of-district contributions meant excluding
donations addressed to national entities (parties and leadership
contestants). As shown in Section~\ref{sec-data-summaries},
contributions of this type represented a majority (around 80\%) of
contributions during our study period. The exclusion of these entities
therefore limits the generalizability of our findings; the majority of
Canadian donation cannot actually be assessed as having been sent in- or
out-of-district. While federal political donation in general is best
understood as a nationalized phenomenon, our study provides evidence
that the individual donor bases for two key players on the Canadian
federal political financing scene -- district associations and electoral
candidates -- are largely regionalized. District associations and
candidates receive some resources through internal party transfers;
however, these entities continue to greatly depend on individual
contributions, making donor geography an important dimension of their
financial support base. Further research may use our dataset of
donor-geolocated political contributions to investigate whether national
party donor bases are similarly constrained to particular regions of the
country.

Third, our chosen study period (October 2015 to April 2024) truncates
the available data in the middle of the 44th parliament. This was done
to restrict our analysis to a single, static set of districts. While we
succeeded in excluding non-2013-order districts, we cannot guarantee
that all valid 2013-order FED entries were captured\footnote{Over all
  entries in the dataset which had recipients belonging to 2013-order
  districts (2004 - present), only 0.5\% were not captured by our chosen
  range.}, since the adoption of new districts is a multi-year
undertaking, and timelines can differ between locations
(\citeproc{ref-canada_redistribution_2025}{Elections Canada 2025b}).
Further, our chosen date range results in the third non-election period
of our analysis being incomplete, so we cannot claim to have studied
three full parliaments. However, we did not identify any systematic
differences in out-of-district donation rates between parliaments,
having mainly discovered variation through space, and not through time.

\newpage
\part*{Appendix}
\appendix

\counterwithin{figure}{section}
\counterwithin{table}{section}
\renewcommand{\thefigure}{\thesection\arabic{figure}}
\renewcommand{\thetable}{\thesection\arabic{table}}

\section{Background on Canadian Political Financing}\label{sec-info}

\subsection{Individual Contribution Limits}\label{sec-info-limits}

The \emph{Canada Elections Act} imposes individual contribution limits
on donors seeking to contribute to one of the following registered,
federal political entities\footnote{A full definition for each entity is
  available at Elections Canada
  (\citeproc{ref-canada_political_2026}{2026d}).}: political parties,
electoral district associations (EDAs), candidates, party nomination
contestants, and party leadership contestants
(\citeproc{ref-canada_annual_2026}{Elections Canada 2026a}). These
limits are shown in Table~\ref{tbl-info-limits-1}\footnote{Contribution
  limits presented here are only for 2026. They increase by \$25 every
  year on January 1st.} and Table~\ref{tbl-info-limits-2}. On top of
these, candidates are also permitted to give an additional \$1,775 in
total per year in contributions to other candidates, registered
associations, or nomination contests, including their own. Regulated
expenses for any of the entities themselves cannot exceed specific
values, which are determined on an annual basis, and are different for
each riding (see \citeproc{ref-canada_expenses_2026}{Elections Canada
2026b}). Individual contributions to registered third-party
organizations are not regulated by the \emph{Act}, and were not included
in this study (\citeproc{ref-canada_understanding_2026}{Elections Canada
2026f}).

Only Canadian citizens or permanent residents can make contributions
(\citeproc{ref-canada_understanding_2026}{Elections Canada 2026f}). Any
donor who contributed \$20 or more to any of the five federal political
entities listed above must have their name recorded (but not reported to
Elections Canada). Donors with contributions exceeding \$200 must also
have their names and addresses recorded and reported to Elections
Canada. Cash donations above \$20 are prohibited
(\citeproc{ref-canada_understanding_2026}{Elections Canada 2026f}).

\begin{table}

\caption{\label{tbl-info-limits-1}2026 Annual Limits on Individual
Contributions, Loans, and Loan Guarantees (adapted from Elections Canada
website).}

\centering{

\centering
\begin{tblr}[         %% tabularray outer open
]                     %% tabularray outer close
{                     %% tabularray inner open
width={1\linewidth},
colspec={X[]X[]},
hline{2}={1-2}{solid, black, 0.05em},
hline{1}={1-2}{solid, black, 0.08em},
hline{6}={1-2}{solid, black, 0.08em},
column{1}={}{halign=l},
column{2}={}{halign=r},
}                     %% tabularray inner close
Political Entity & 2026 Annual Limit (CAD) \\
To each registered national party & \$1,775 \\
In total, to all leadership contestants in a particular contest & \$1,775 \\
In total, to all EDAs, nomination contestants, and candidates (for each registered party) & \$1,775 \\
In total, to each independent or non-affiliated candidate & \$1,775 \\
\end{tblr}

}

\end{table}%

\begin{table}

\caption{\label{tbl-info-limits-2}2026 Limits on Individual
Contributions, Loans, and Loan Guarantees, for Politicians Contributing
to Their Own Campaigns.}

\centering{

\centering
\begin{tblr}[         %% tabularray outer open
]                     %% tabularray outer close
{                     %% tabularray inner open
width={1\linewidth},
colspec={X[]X[]},
hline{2}={1-2}{solid, black, 0.05em},
hline{1}={1-2}{solid, black, 0.08em},
hline{5}={1-2}{solid, black, 0.08em},
column{1}={}{halign=l},
column{2}={}{halign=r},
}                     %% tabularray inner close
Politician & Own-Campaign Contribution Limit (CAD) \\
Nomination Contestant & \$1,000 (in addition to regular annual limits) \\
Candidate & \$5,000 \\
Leadership Contestant & \$25,000 \\
\end{tblr}

}

\end{table}%

\subsection{Financial Reporting Requirements of Federal Political
Entities}\label{sec-info-reporting}

All of the information presented in this section can be found on the
``Political Financing'' page of the Elections Canada website
(\citeproc{ref-canada_financing_2026}{Elections Canada 2026e}).
Financial returns represent an account of an entity's inflow
(i.e.~contributions, loans or transfers) and outflow (expenses). At
minimum, all entities must submit financial returns annually (for
parties and EDAs), or within 4 to 8 months after an electoral event (for
parties, candidates, third parties, leadership contestants, and
nomination contestants exceeding \$1000 in contributions or expenses).
Returns with large transaction amounts must be accompanied by an
external audit. Entities such as parties and EDAs may be required to
file reports on a quarterly basis. Some (but not all) returns are
reviewed by Elections Canada auditors. Risk-based auditing is used to
determine which files will be audited and what level of audit will be
performed (\citeproc{ref-canada_audit_2026}{Elections Canada 2026c}).

\subsection{Electoral Redistricting}\label{sec-info-redistricting}

Electoral district boundaries are subject to change every 10 years, as
required by the \emph{Constitution of Canada}
(\citeproc{ref-canada_redistribution_2025}{Elections Canada 2025b}). The
2013 representational order consisted of 338 districts, and the 2023
representational order consisted of 343 districts. The boundaries of
existing districts are subject to change as well. The 2013
representational order was decidedly in effect between the calling of
the 42nd general election on October 2, 2015, and the first dissolution
of Parliament occurring after April 22, 2024
(\citeproc{ref-canada_redistribution_effect_2023}{Elections Canada
2023}).

\section{Donations Data Pipeline}\label{sec-pipeline}

In this section we describe the three major cleaning and validation
steps required to transform contributions data from Elections Canada
into the dataset used throughout our analysis. These scripts rely on
other government datasets, cleaned separately, and documented in
Appendix~\ref{sec-data-supplements}.

\subsection{Investigative Journalism Foundation
Pre-Cleaning}\label{sec-source-to-IJF}

The original contribution datasets used in this study can be found on
Elections Canada's ``Open data on contributions'' page, under the
heading ``Data sets updated'' (\citeproc{ref-canada_open_2025}{Elections
Canada 2025a}). Two products are offered. The first is a dataset of
contribution reports filed by all Canadian registered federal political
entities from January 2004 to present
(\citeproc{ref-canada_open_2025}{Elections Canada 2025a}). The second is
an audited version of the first. Due to pre-screening policies and
administrative lag, only some of the first dataset is audited. In our
case, the audited version of the dataset had only 61.7\% as many entries
as the full un-audited data (n = 10095198; both accessed on May 5,
2026). Both datasets are updated on a weekly basis by Elections Canada
(\citeproc{ref-canada_open_2025}{Elections Canada 2025a}), as new audits
and financial reports are filed.

Staff at the Investigative Journalism Foundation created a data pipeline
to aggregate, de-duplicate, and clean these datasets. Both audited and
un-audited versions of the same record were available for most records.
To avoid duplicating these, audited and un-audited versions of the same
record were linked by assuming that rows sharing the same donation date,
electoral event, recipient, and recipient political party referred to
the same entry. We note that this linkage approach carries the small
risk that two distinct donations to the same recipient on the same date
are incorrectly treated as duplicates. The pipeline preferentially keeps
audited records when they are available.

Further de-duplication was effectuated differently based on the type of
contribution being recorded. For loans, sums reported during an Annual
filing might appear in a Quarterly report filed in the same year. In
such cases, Annual records were preserved and Quarterly records were
discarded. For all other types of contributions, only true duplicates
(i.e.~rows with identical entries identical in every field) were
discarded. Additionally, some entries represented aggregated sums of
contributions from anonymous donors. For these, Elections Canada reports
a cumulative total per recipient Annually and Quarterly, meaning that
the same aggregate sum could appear twice in the same year. To avoid
double-counting these, only the largest reported aggregate sum per
recipient per year was retained. Finally, a few smaller cleaning
procedures were applied to the columns themselves;
Table~\ref{tbl-source-to-IJF-1} shows where each of the columns present
in the IJF data were drawn from columns at the source, and any notable
procedures applied to them.

\begin{table}

\caption{\label{tbl-source-to-IJF-1}Description of where the columns in
the IJF data were derived from the source, and any cleaning procedures
applied. Variables added for internal database consistency
(e.g.~processing dates) and redundant variables (e.g.~donation year) are
not shown here.}

\centering{

\centering
\begin{tblr}[         %% tabularray outer open
]                     %% tabularray outer close
{                     %% tabularray inner open
width={1\linewidth},
colspec={X[0.2]X[0.4]X[0.4]},
hline{2}={1-3}{solid, black, 0.05em},
hline{1}={1-3}{solid, black, 0.08em},
hline{13}={1-3}{solid, black, 0.08em},
column{1-3}={}{halign=l},
}                     %% tabularray inner close
IJF Column & EC Source Column(s) & Notes \\
donor\_full\_name & Contributor name ; Contributor last name ; Contributor first name ; Contributor middle initial & Concatenated source columns, delimited by commas. \\
donor\_type & Contributor type & No changes from source. \\
donor\_location & Contributor City ; Contributor Province ; Contributor Postal code & Concatenated source columns, delimited by commas. \\
donation\_date & Contribution Received date ; Fiscal/Election date & Used `Contribution Received date` when available, filled with Fiscal/Election date otherwise. Flagged when imputation occured. Removed invalid dates. \\
electoral\_event & Electoral event & Donations to leadership and nomination contestants had NA entries here. Otherwise, no changes from source. \\
amount\_monetary & Monetary amount & No changes from source. \\
amount\_non\_monetary & Non-Monetary amount & No changes from source. \\
recipient & Recipient ; Recipient last name ; Recipient first name ; Recipient middle initial & Concatenated source columns, delimited by commas. \\
electoral\_district & Electoral District & No changes from source. \\
political\_entity & Political Entity & No changes from source. \\
political\_party & Political Party of Recipient & No changes from source. \\
\end{tblr}

}

\end{table}%

\subsection{Initial Cleaning of IJF Data}\label{sec-cleaning-IJF}

We applied further cleaning and filtering steps to the data supplied by
the Investigative Journalism Foundation. First, we removed redundant
columns (e.g.~data had columns \texttt{year} and \texttt{donation\_year}
which were identical for every entry) and columns used for internal
consistency (e.g.~the date the entry was added by the organization).
Next, we removed entries with malformatted dates, or disagreements
between the year and date columns. We then filtered for only donations
stated to have been made by individuals. Roughly 0.4\% of the data was
removed at this step, the overwhelming majority of which were records
dated before 2006, outside the scope of our study. Next, we filtered for
donations where the sum of monetary and non-monetary contributions was
present and positive. We then removed entries where the donor or
recipient name was either missing or shorter than 4 characters (after
trimming whitespaces). Finally, we removed individual donations with
amounts larger than \$25000. Here we were careful to exclude entries
that represented aggregated sums of anonymous contributions, which were
flagged by searching for entries where \texttt{donor\_full\_name}
started with ``Contribut''. In total, roughly 0.41\% of entries (n
\textasciitilde{} 30,000) were discarded in this section.

Donor postal codes were extracted from the \texttt{donor\_location}
column using regex-based string matching. Canadian postal codes always
follow the format A0A0A0, where 0 is a digit and A is a letter
(excluding I, O, D, F, Q, or U). Also, they may not start with W or Z
(\citeproc{ref-canada_post_guidelines_2026}{Canada Post 2026}). These
facts were used to verify the structure of extracted postal codes, and
apply the following imputations:

\begin{itemize}
\tightlist
\item
  ``O''s in a number slot were taken to mean ``0''
\item
  ``I''s in a number slot were taken to mean ``1''
\item
  ``!'' in a number slot were taken to mean ``1''
\end{itemize}

After imputations, we were able to recover 99.89\% of individual donor
postal codes. We note that while these postal codes are guaranteed to be
structurally correct, we did not validate whether the codes matched
their associated provinces, cities, or were actually used. The effective
validity of these postal codes is verified in
Appendix~\ref{sec-IJF-to-donations-data}. 0.35\% of entries (n
\textasciitilde{} 26,000) represented aggregated sums of anonymous
contributions, and therefore did not have associated postal codes.
Despite their scarcity, these amounted to 24.1\% of the total
contributed amount in the dataset. These were kept at this stage.

We also performed simple string cleaning procedures such as whitespace
stripping and entity matching (e.g.~United Party of Canada (UP) and
United Party of Canada clearly refer to the same political party).
Notably all donation recipients in the dataset had associated electoral
districts (apart from leadership contestants and registered parties,
since these are non-local entities) and political parties. Spot checking
revealed no district or party attribution errors among a small set of
recipients, however we could not completely confirm that recipient
districts or parties were completely correct.

\subsection{\texorpdfstring{Clean IJF to
\texttt{donations\_data}}{Clean IJF to donations\_data}}\label{sec-IJF-to-donations-data}

In this section we describe how our primary study dataset was derived
from the cleaned IJF data, and augmented to include the donor's district
using a custom postal code conversion file (PCCF) described in
Appendix~\ref{sec-data-supplements-spatial}. The resulting dataset
represents all donor-geolocated contributions above \$200 made out to
federal political entities (excluding nomination contestants and
candidates during by-elections).

First, we restricted the dataset to only those entries with donation
dates between October 2, 2015, and April 22, 2024, capturing the
effective period of the 2013 representational order. Next, we removed
entries where the recipient's district was not one of the 338 possible
districts in the 2013 representational order, referencing our modified
PCCF. Next, we removed entries where the reported electoral event did
not happen within our date range (e.g.~39th general election was in
2006). Although it is possible that these were retroactively added for
past events, less than 0.001\% of the data had this issue, so their
removal did not impact our analysis. Next, we removed individual
donations addressed to localized entities which had amounts larger than
what the contribution limits allowed. For candidates this was \$5000,
and for registered associations this was \$3550. Next, we removed
donations addressed to leadership contestants for parties where no such
contest occurred during our study period. Nomination contestants were
removed entirely, to simplify analysis. 3\% of donations (7\% of summed
total) received during general election periods had receipt dates
outside the stated period (usually dated after the writ period had
ended, not before it started). These were removed when fitting our
regression models, but kept in most descriptive visulaizations.

As mentioned in the main text, our study considers only donations above
\$200. This discards 45\% of the total contributed amount during our
study period, but only 32\% of the total amount received by localized
entities (i.e.~candidates and district associations). These excluded
data are discussed further in Appendix~\ref{sec-missing-data}. For the
remaining entries, donor postal codes were matched to their
corresponding census dissemination areas (DAs) and federal electoral
districts (FEDs) using our modified PCCF file. 1.2\% of donor postal
codes (representing 1.3\% of the total amount contributed) were either
missing or could not be matched, and were dropped (2.1\% for localized
entities). Figure~\ref{fig-IJF-to-donations-data-missingness} breaks
down the discarded data by its cause, for localized entities. Finally,
columns were renamed, merged, and variables were joined to improve the
dataset's usability. The resulting schema is shown in part in
Table~\ref{tbl-IJF-to-donations-data-schema}.

\begin{table}

\caption{\label{tbl-IJF-to-donations-data-schema}Primary variables used
in our study, with examples and descriptions. Although not shown,
additional information used such as census variables or geographic
attributes of districts are easily retrived by relating a donor's
dissemination area or district to custom-made lookup tables.}

\centering{

\centering
\begin{tblr}[         %% tabularray outer open
]                     %% tabularray outer close
{                     %% tabularray inner open
width={1\linewidth},
colspec={X[]X[]X[]},
hline{2}={1-3}{solid, black, 0.05em},
hline{1}={1-3}{solid, black, 0.08em},
hline{13}={1-3}{solid, black, 0.08em},
cell{1}{1}={}{font=\bfseries\fontsize{0.85em}{1.15em}\selectfont, halign=l},
cell{1}{2}={}{font=\bfseries\fontsize{0.85em}{1.15em}\selectfont, halign=r},
cell{1}{3}={}{font=\bfseries\fontsize{0.85em}{1.15em}\selectfont, halign=r},
cell{2-12}{1}={}{font=\fontsize{0.85em}{1.15em}\selectfont, halign=l},
cell{2-12}{2}={}{font=\fontsize{0.85em}{1.15em}\selectfont, halign=r},
cell{2-12}{3}={}{font=\fontsize{0.85em}{1.15em}\selectfont, halign=r},
}                     %% tabularray inner close
Variable & Example & Description \\
donor\_name & Alexandra Lulka & Donor full name. \\
donor\_dissemination\_area & 35202387 & Dissemination area (UID) associated with donor postal code. \\
donor\_district & Eglinton--Lawrence & Electoral district associated with donor postal code. \\
donation\_date & 2015-10-18 & Date donation was received. \\
electoral\_event & General Election & Electoral event (if any) occurring during donation. \\
total\_amount & 563.3 & Summed monetary and non-monetary value of the donation, in CAD. \\
recipient\_name & Joe Oliver & Recipient full name. \\
recipient\_district & Eglinton--Lawrence & District of the recipient (if applicable). \\
political\_entity & Candidates & Entity type of the recipient. \\
political\_party & Conservative Party of Canada & Recipient party affiliation (if any). \\
is\_out\_of\_district & FALSE & Whether donor and recipient districts are different. \\
\end{tblr}

}

\end{table}%

\begin{figure}

\centering{

\pandocbounded{\includegraphics[keepaspectratio]{paper_files/figure-pdf/fig-IJF-to-donations-data-missingness-1.pdf}}

}

\caption{\label{fig-IJF-to-donations-data-missingness}Total amount
donated to localized entities (candidates and district associations)
during our study period, split by whether or not the share was included
in the study. The share of discarded contributions (left) is further
broken down by the cause of its removal. We also display the share of
included data to which not all census variables could be attributed
(light blue).}

\end{figure}%

\section{Supplementary Datasets}\label{sec-data-supplements}

In this section we describe how we created two custom lookup tables that
the cleaning scripts in Appendix~\ref{sec-pipeline} rely on; one to
convert between different levels of spatial hierarchy, and a second to
retrieve specific socio-demographic statistics at the census
dissemination area. Their place in the full pipeline is illustrated in
\textbf{?@fig-data-pipeline}.

\subsection{Spatial Lookup Table}\label{sec-data-supplements-spatial}

The purpose of the table constructed in this section is to provide
unambiguous links between any administrative area of interest and
another at a larger scale (e.g.~postal code to district, dissemination
area to province, etc.). This custom table is a version of Statistics
Canada's 2024 postal code conversion files (PCCF), modified to resolve
many-to-many spatial overlaps between postal codes, dissemination areas
(DAs), and electoral districts (FEDs). The PCCF we used provides
linkages between all registered Canadian postal codes (accurate as of
June 2024) and various administrative divisions, including the
dissemination areas (DAs) of the 2021 Census, and the federal electoral
districts (FEDs) of the 2013 representational order
(\citeproc{ref-statcan_pccf_2025}{Statistics Canada 2025}). These sets
of DAs and FEDs are the same versions we chose for our study. Access to
this PCCF was granted to us by the University of Toronto. The result of
this process is shown in part in Table~\ref{tbl-spatial-final}.

To resolve cases where one postal code straddled multiple dissemination
areas, we used a single link indicator (SLI) to establish a one-to-one
relationship between postal codes and DAs, available in PCCF files by
default (\citeproc{ref-statcan_pccf_2025}{Statistics Canada 2025}).
Statistics Canada determines SLIs by estimating the population that
resides in each DA within the postal code area (based on 2021 Census
population counts) and assigns the code to the one with a higher
population (\citeproc{ref-statcan_pccf_2025}{Statistics Canada 2025}).
This procedure removed 35\% of links present in the original PCCF,
however no single postal code was removed from the file entirely. A
small number of postal codes (0.7\%) were removed for missing either an
associated DA or FED. All postal codes had associated provinces. In
addition, 1.6\% of dissemination areas in the dataset straddled multiple
district boundaries. To resolve this, for each overlap case, we (1)
counted how many postal codes within the dissemination area linked to
each candidate FED, and (2) assigned the entire DA, and all of its
postal codes to the FED holding the plurality of those postal codes,
discarding the postal codes that mapped to any minority FED. Ties
between FEDs were broken deterministically by FED UID. This resolution
step resulted in the loss of roughly 0.5\% of postal codes. Applied to
the contributions dataset used in this study, these resolutions
increased the share of entries whose donor postal codes could not be
matched to the PCCF from 0.7\% to 1.2\%.

As a final step, we ran a suite of tests, primarily to confirm that the
resolved table formed a strictly nested hierarchy: each postal code
mapped to a single DA, each DA or postal code to a single FED, and each
FED, DA, or postal code to a single province. We also verified that
every postal code appeared exactly once, and that they were correctly
formatted following Canada Post's guidelines
(\citeproc{ref-canada_post_guidelines_2026}{Canada Post 2026}). We also
confirmed that all DA and FED UIDs could be cross-referenced with
Statistics Canada's 2021 DA and 2013 FED shapefiles
(\citeproc{ref-statcan_boundary_2021}{Statistics Canada 2022a}), and
that all 338 districts, 10 provinces, and 3 territories remained
represented. Because of these tests, this spatial lookup table was
treated as an authoritative list of valid postal codes, dissemination
areas, districts, and provinces, which greatly simplified our analysis.

\begin{table}

\caption{\label{tbl-spatial-final}Custom spatial lookup table, modified
from Statistics Canada's Postal Code Conversion File. This table
unambiguously links each Canadian postal code to a 2021 census
dissemination area, and each census dissemination area to a federal
electoral district from the 2013 representational order. Only a sample
of four entries are shown.}

\centering{

\centering
\begin{tblr}[         %% tabularray outer open
]                     %% tabularray outer close
{                     %% tabularray inner open
width={1\linewidth},
colspec={X[]X[]X[]},
hline{2}={1-3}{solid, black, 0.05em},
hline{1}={1-3}{solid, black, 0.08em},
hline{6}={1-3}{solid, black, 0.08em},
cell{1}{1}={}{font=\bfseries\fontsize{0.85em}{1.15em}\selectfont, halign=l},
cell{1}{2}={}{font=\bfseries\fontsize{0.85em}{1.15em}\selectfont, halign=r},
cell{1}{3}={}{font=\bfseries\fontsize{0.85em}{1.15em}\selectfont, halign=r},
cell{2-5}{1}={}{font=\fontsize{0.85em}{1.15em}\selectfont, halign=l},
cell{2-5}{2}={}{font=\fontsize{0.85em}{1.15em}\selectfont, halign=r},
cell{2-5}{3}={}{font=\fontsize{0.85em}{1.15em}\selectfont, halign=r},
}                     %% tabularray inner close
Postal Code & Dissemination Area & Federal Electoral District \\
M4P1E4 & 35202387 & Eglinton--Lawrence \\
K1A0A6 & 35061205 & Ottawa Centre \\
V6B1A1 & 59153201 & Vancouver Centre \\
H3A0G4 & 24661048 & Outremont \\
\end{tblr}

}

\end{table}%

\subsection{Census Lookup Table}\label{sec-data-supplements-census}

The purpose of the table constructed in this section is to provide
specific socio-demographic variables (listed in
Table~\ref{tbl-data-supplements-census}) for every dissemination area
(DA) present in the spatial lookup table (described in
Section~\ref{sec-data-supplements-spatial}). Socio-demographic
information at this spatial scale was sourced from the 2021 Census, and
specifically from the ``Canada, provinces, territories, census divisions
(CDs), census subdivisions (CSDs) and dissemination areas (DAs)''
available for download online
(\citeproc{ref-statcan_census_zip_2022}{Statistics Canada 2022b}). We
chose the 2021 Census over the 2016 Census since the midpoint of our
time period (2019) is closer to 2021 than 2016. We chose the
dissemination area, an administrative region that houses roughly 400 to
700 people, since it is the smallest spatial scale where each of our
desired statistics are recorded
(\citeproc{ref-statcan_da_2021}{Statistics Canada 2021}). All DAs listed
in our spatial lookup file (see
Appendix~\ref{sec-data-supplements-spatial}) were present in the Census
dataset we used. Since these DAs are the only ones we would ever use for
analysis, the other 15.1\% of DAs were dropped.

3.2\% of remaining DAs were missing at least one census variable: 0.05\%
were missing all demographic variables, 0.43\% only had population,
area, and density, 0.12\% were only missing the 25\% sampled variables
(i.e.~ethnicity and education level), and 2.62\% had all variables
recorded except household income. Notably, all entries' missingness
could be grouped into five strictly hierarchical stages, summarized in
Table~\ref{tbl-census-missingness}. Entries with missing variables were
not dropped, but were instead flagged so that they could be easily
identified during analysis. We also converted ethnicity and education
counts to proportions. Table~\ref{tbl-census-final} provides a glimpse
of the resulting Census dataset.

\begin{table}

\caption{\label{tbl-data-supplements-census}Census IDs, descriptions,
and sample sizes for each census variable used in this study. Variables
with a 25\% sample size were recorded for only 25\% of the population,
to preserve privacy.}

\centering{

\centering
\begin{tblr}[         %% tabularray outer open
]                     %% tabularray outer close
{                     %% tabularray inner open
width={1\linewidth},
colspec={X[]X[]X[]},
hline{2}={1-3}{solid, black, 0.05em},
hline{1}={1-3}{solid, black, 0.08em},
hline{21}={1-3}{solid, black, 0.08em},
column{1}={}{halign=l},
column{2-3}={}{halign=r},
}                     %% tabularray inner close
Characteristic ID & Description & Sample Size \\
1 & Population, 2021 & 100\% \\
6 & Population density (per square kilometres) & 100\% \\
7 & Land area (square kilometres) & 100\% \\
39 & Average (mean) age & 100\% \\
243 & Median household income (total, in 2020) & 100\% \\
1683 & Total population in private households & 25\% \\
1684 & Total visible minority population & 25\% \\
1685 & Total South Asian population & 25\% \\
1686 & Total Chinese population & 25\% \\
1687 & Total Black population & 25\% \\
1688 & Total Filipino population & 25\% \\
1689 & Total Arab population & 25\% \\
1690 & Total Latin American population & 25\% \\
1691 & Total Southeast Asian population & 25\% \\
1692 & Total West Asian population & 25\% \\
1693 & Total Korean population & 25\% \\
1694 & Total Japanese population & 25\% \\
1998 & Total population aged 15+ & 25\% \\
2001 & Total population with a postsecondary certificate & 25\% \\
\end{tblr}

}

\end{table}%

\begin{table}

\caption{\label{tbl-census-missingness}Structure of missing Census data.
Missingness is hierarchical in the Census dataset: entries at a higher
stage of missingness are always missing everything in the lower stages.}

\centering{

\centering
\begin{tblr}[         %% tabularray outer open
]                     %% tabularray outer close
{                     %% tabularray inner open
width={1\linewidth},
colspec={X[]X[]X[]},
hline{2}={1-3}{solid, black, 0.05em},
hline{1}={1-3}{solid, black, 0.08em},
hline{7}={1-3}{solid, black, 0.08em},
column{1}={}{halign=l},
column{2-3}={}{halign=r},
}                     %% tabularray inner close
Missing Stage & Count & \% of Total \\
0. Complete & 47593 & 96.78\% \\
1. Missing Median Household Income & 1288 & 2.62\% \\
2. Missing 25\% Sample Variables & 57 & 0.12\% \\
3. Missing Age & 211 & 0.43\% \\
4. Missing Population & 26 & 0.05\% \\
\end{tblr}

}

\end{table}%

\begin{table}

\caption{\label{tbl-census-final}Custom census lookup table. This table
provides desired socio-demographic information for every dissemination
area in the 2021 Census. Only a sample of four entries are shown. Only
population, density, and median household income are shown.}

\centering{

\centering
\begin{tblr}[         %% tabularray outer open
]                     %% tabularray outer close
{                     %% tabularray inner open
width={1\linewidth},
colspec={X[]X[]X[]X[]},
hline{2}={1-4}{solid, black, 0.05em},
hline{1}={1-4}{solid, black, 0.08em},
hline{6}={1-4}{solid, black, 0.08em},
cell{1}{1}={}{font=\bfseries\fontsize{0.85em}{1.15em}\selectfont, halign=l},
cell{1}{2}={}{font=\bfseries\fontsize{0.85em}{1.15em}\selectfont, halign=r},
cell{1}{3}={}{font=\bfseries\fontsize{0.85em}{1.15em}\selectfont, halign=r},
cell{1}{4}={}{font=\bfseries\fontsize{0.85em}{1.15em}\selectfont},
cell{2-5}{1}={}{font=\fontsize{0.85em}{1.15em}\selectfont, halign=l},
cell{2-5}{2}={}{font=\fontsize{0.85em}{1.15em}\selectfont, halign=r},
cell{2-5}{3}={}{font=\fontsize{0.85em}{1.15em}\selectfont, halign=r},
cell{2-5}{4}={}{font=\fontsize{0.85em}{1.15em}\selectfont},
}                     %% tabularray inner close
Dissemination Area (UID) & Population & Density (per sqkm) & Median Household Income (CAD) \\
35202387 & 612 & 4381.2 & \$98,500 \\
35061205 & 548 & 3920.7 & \$87,200 \\
59153201 & 701 & 6210.5 & \$76,300 \\
24661048 & 489 & 5104.9 & \$92,100 \\
\end{tblr}

}

\end{table}%

\section{Characterizing Excluded Contributions Under
\$200}\label{sec-missing-data}

In our analysis of out-of-district political donations addressed to
electoral candidates and district associations, we excluded donations
under \$200, and donations whose donors could not successfully be
geolocated, which in total represents 33.9\% of the total amount
received by these entities during our study period (see
Table~\ref{tbl-missing-data-cases}). In this section, we compare the
recipient characteristics of contributions under \$200 with those kept
in our study to assess the impact of their removal on our analysis. By
the \emph{Canada Elections Act}, donors during a given calendar year (or
writ period for donations to candidates) who contributed less than \$200
in total are not required to have their name or address disclosed to
Elections Canada (\citeproc{ref-canada_understanding_2026}{Elections
Canada 2026f}). Donations whose donors did not disclose this personal
information are recorded in our dataset as aggregated sums, grouped by
recipient and time-period (writ period for candidates, and quarterly or
annual reports for district associations). Otherwise, donations under
\$200 appear as individual entries with associated donor names and
locations. These two cases are compared against the data separately, and
shown in Figure~\ref{fig-missing-data-distribution-anonymous}, and
Figure~\ref{fig-missing-data-distribution-disclosed} respectively.

Anonymous contributions under \$200 are the largest source of missing
data, and are overall similar in distribution to our study data across
the factors considered
(Figure~\ref{fig-missing-data-distribution-anonymous}). A larger
proportion of these were addressed to district associations, Liberal
entities over Conservative entities, and earlier in the decade (2015 to
2018, and 42nd general election) rather than later (2021 to 2023, and
44th general election). Anonymous contributions were nearly identically
distributed to those used in our study with respect to the region of the
receiving entity. No single factor among those considered in the
anonymous contributions deviated from the study data by more than 10\%.

Donor-geolocated contributions under \$200 were distributed differently
than the data in our study
(Figure~\ref{fig-missing-data-distribution-disclosed}). A majority
(83\%) of these contributions were addressed to Liberal district
associations specifically, largely outside general election periods (see
Panel E in Figure~\ref{fig-missing-data-distribution-disclosed}). These
contributions are similarly distributed to our study data across
regions, and in time for those few candidate-bound cases. It is unclear
why donors who donate to Liberal district associations are more likely
to have their donor locations disclosed.

We also found that only 18\% of these donations were sent
out-of-district (21.5\% of the total amount contributed), compared to
42\% and 46\% in the studied dataset. While we cannot confirm these
numbers for the anonymous contributions, this suggests that
contributions under \$200 are overall less likely to leave their
districts of origin. This is also supported by the positive association
between contributed amount and out-of-district propensity found in our
study. While this suggests that our point estimates of out-of-district
donation rates are overestimates of the true out-of-district propensity
(especially the amount-unweighted proportions), the similarity in
distributions between these and contributions above \$200 (especially
across regions) suggests that excluding these contributions did not
greatly impact our model findings.

\begin{table}

\caption{\label{tbl-missing-data-cases}Breakdown of all contributions
addressed to district associations and candidates (during general
election periods) from 2015 to 2024, by cause of exclusion from
analysis. Study Data is included for reference. We display the share of
the total amount received by these entities capture by each case.}

\centering{

\centering
\begin{tblr}[         %% tabularray outer open
]                     %% tabularray outer close
{                     %% tabularray inner open
width={1\linewidth},
colspec={X[]X[]},
hline{2}={1-2}{solid, black, 0.05em},
hline{1}={1-2}{solid, black, 0.08em},
hline{6}={1-2}{solid, black, 0.08em},
cell{1}{1}={}{font=\bfseries\fontsize{0.85em}{1.15em}\selectfont, halign=l},
cell{1}{2}={}{font=\bfseries\fontsize{0.85em}{1.15em}\selectfont, halign=r},
cell{2-5}{1}={}{font=\fontsize{0.85em}{1.15em}\selectfont, halign=l},
cell{2-5}{2}={}{font=\fontsize{0.85em}{1.15em}\selectfont, halign=r},
}                     %% tabularray inner close
Data Case & Share of Total Contribution Amount \\
Study Data & 66.1\% \\
Unsuccessfully Donor-Geolocated (Over \$200) & 1.6\% \\
Donor-Geolocated (Under \$200) & 7.9\% \\
Anonymous (Under \$200) & 24.4\% \\
\end{tblr}

}

\end{table}%

\begin{figure}

\centering{

\pandocbounded{\includegraphics[keepaspectratio]{paper_files/figure-pdf/fig-missing-data-distribution-anonymous-1.pdf}}

}

\caption{\label{fig-missing-data-distribution-anonymous}Comparing the
distributions of donations to candidates and district associations above
\$200 whose donors were successfully situated in their districts (blue;
those included in our study) vs.~anonymous donations under \$200
(orange; not included in the study). Datasets are compared across five
factors, and distributions are weighted by total amount contributed.
Panels D and E only compare datasets for donations bound for candidates
and district associations respectively.}

\end{figure}%

\begin{figure}

\centering{

\pandocbounded{\includegraphics[keepaspectratio]{paper_files/figure-pdf/fig-missing-data-distribution-disclosed-1.pdf}}

}

\caption{\label{fig-missing-data-distribution-disclosed}Comparing the
distributions of donations to candidates and district associations above
\$200 whose donors were successfully situated in their districts (blue;
those included in our study) vs.~location-disclosed donations under
\$200 (yellow; not included in the study). Datasets are compared across
five factors, and distributions are weighted by total amount
contributed. Panels D and E only compare datasets for donations bound
for candidates and district associations respectively.}

\end{figure}%

\section{Alternate Model
Specifications}\label{sec-alternate-model-specifications}

Here we present the results of six alternate out-of-district propensity
models, mainly fit to assess the robustness of our results in
Chapter~\ref{sec-results} to other resonable model specifications.

Table~\ref{tbl-results-model-individual-EDA} shows results from an
individual-level model with the same covariates, fit using only
donations bound for district associations (i.e.~excluding those received
by candidates). We observe similar coefficients and overall model fit.
The only major difference between this and the primary model is that
average age is negatively associated with out-of-district propensity,
with statistical significance at the \(\alpha = 0.05\) level.
Table~\ref{tbl-results-model-district-EDA} shows results from a
district-level model with the same covariates, fit using only donations
bound for district associations (i.e.~excluding those received by
candidates). We observe similar coefficients and overall model fit.

Table~\ref{tbl-results-model-individual-candidiate} shows results from
an individual-level model with the same covariates, fit using only
donations received by candidates (i.e.~excluding those received by
district associations). We observe similar coefficient and overall model
fit. Table~\ref{tbl-results-model-individual-candidiate-competitive}
shows results from a slighly more elaborate version of this model, that
includes the election period in which the donation was received, as well
as how competitive the donor's district was during that period. Both of
these more specific covariates were significantly assocated with
out-of-district propensity, but did not affect the rough magnitudes and
significances anywhere else in the model.

Table~\ref{tbl-results-model-individual-disaggregated} and
Table~\ref{tbl-results-model-district-disaggregated} shows results from
individual level and district level models respectively, where province
and visible minority proportion were disaggregated into more specific
categories. In both models, the coefficients associated with provinces
each behaved similarly to their parent region. The proportion of
habitants belonging to all ethnic groups available in the Census were
positively associated with out-of-district donation at the individual
level. However, only the proportion of South Asian, Black, and Chinese
inhabitants of a district was associated with it's share of donations
that were bound for other districts. Further, disaggregation greatly
increased the number of ocefficients used to fit the model, for marginal
gains in the model's explanatory power. For example, the disaggregated
district level model has over double the number of coefficients compared
to the one included in the main text, for only a single percentage
increase in the model's explainability. For this reason, grouped
variables were kept in the main models at both levels.

\begin{table}

\caption{\label{tbl-results-model-individual-EDA}Results from a logistic
regression model where the outcome is whether or not an individual
donation was sent out-of-district. Only donations sent to district
assocaitions are considered. The reference party is the Conservative
Party, and the reference region is Ontario. Estimates are bolded when
significant at the alpha = 0.05 level.}

\centering{

\centering
\begin{tblr}[         %% tabularray outer open
]                     %% tabularray outer close
{                     %% tabularray inner open
colspec={Q[]Q[]Q[]Q[]Q[]},
hline{2}={5}{solid, black, 0.03em},
hline{2}={2,4}{solid, black, 0.03em, l=-0.5},
hline{2}={3}{solid, black, 0.03em, r=-0.5},
hline{3}={1-5}{solid, black, 0.05em},
hline{21}={1-5}{solid, black, 0.05em},
hline{1}={1-5}{solid, black, 0.08em},
hline{27}={1-5}{solid, black, 0.08em},
cell{1}{1}={}{halign=l},
cell{1}{2}={c=2}{halign=c},
cell{1}{3}={}{halign=c},
cell{1}{4}={c=2}{halign=c},
cell{1}{5}={}{halign=c},
cell{2-3,5,7-8,10-20}{2}={}{font=\bfseries\fontsize{0.85em}{1.15em}\selectfont, halign=c},
cell{2-3,5,7-8,10-20}{3}={}{font=\bfseries\fontsize{0.85em}{1.15em}\selectfont, halign=c},
cell{2-3,5,7-8,10-20}{4}={}{font=\bfseries\fontsize{0.85em}{1.15em}\selectfont, halign=c},
cell{2-3,5,7-8,10-20}{5}={}{font=\bfseries\fontsize{0.85em}{1.15em}\selectfont, halign=c},
cell{2}{1}={}{font=\bfseries\fontsize{0.85em}{1.15em}\selectfont, halign=l},
cell{3-26}{1}={}{font=\fontsize{0.85em}{1.15em}\selectfont, halign=l},
cell{4,6,9,21-26}{2}={}{font=\fontsize{0.85em}{1.15em}\selectfont, halign=c},
cell{4,6,9,21-26}{3}={}{font=\fontsize{0.85em}{1.15em}\selectfont, halign=c},
cell{4,6,9,21-26}{4}={}{font=\fontsize{0.85em}{1.15em}\selectfont, halign=c},
cell{4,6,9,21-26}{5}={}{font=\fontsize{0.85em}{1.15em}\selectfont, halign=c},
}                     %% tabularray inner close
& Coefficient &  & Odds Ratio &  \\
& Est. & S.E. & Est. & S.E. \\
(Intercept) & \num{-2.789} & \num{0.069} & \num{0.062} & \num{0.004} \\
Conservative & - &  & - &  \\
Liberal & \num{0.367} & \num{0.016} & \num{1.443} & \num{0.023} \\
N.D.P & \num{0.020} & \num{0.022} & \num{1.021} & \num{0.022} \\
Other & \num{0.153} & \num{0.038} & \num{1.166} & \num{0.044} \\
Amount & \num{0.000} & \num{0.000} & \num{1.000} & \num{0.000} \\
Ontario & - &  & - &  \\
Quebec & \num{0.294} & \num{0.023} & \num{1.342} & \num{0.030} \\
British Columbia & \num{-0.427} & \num{0.019} & \num{0.653} & \num{0.013} \\
Prairies & \num{-0.419} & \num{0.021} & \num{0.658} & \num{0.014} \\
Atlantic Canada & \num{-0.254} & \num{0.032} & \num{0.776} & \num{0.025} \\
Territories & \num{-1.842} & \num{0.152} & \num{0.159} & \num{0.024} \\
Age & \num{-0.004} & \num{0.001} & \num{0.996} & \num{0.001} \\
Income & \num{0.000} & \num{0.000} & \num{1.000} & \num{0.000} \\
Education & \num{1.890} & \num{0.066} & \num{6.617} & \num{0.438} \\
Visible Minority & \num{2.054} & \num{0.032} & \num{7.801} & \num{0.252} \\
Margin of Victory (mean) & \num{0.008} & \num{0.001} & \num{1.008} & \num{0.001} \\
Population Density (sqrt) & \num{0.003} & \num{0.000} & \num{1.003} & \num{0.000} \\
Num.Obs. & \num{107442} &  & \num{107442} &  \\
AIC & \num{130829.1} &  & \num{130829.1} &  \\
BIC & \num{130982.5} &  & \num{130982.5} &  \\
Log.Lik. & \num{-65398.556} &  & \num{-65398.556} &  \\
F & \num{914.461} &  & \num{914.461} &  \\
RMSE & \num{0.46} &  & \num{0.46} &  \\
\end{tblr}

}

\end{table}%

\begin{table}

\caption{\label{tbl-results-model-district-EDA}Results from a beta
regression model where the outcome is the out-of-district share of the
sum of EDA-bound contributions over \$200 originating from an individual
district during our study period. The reference region is Ontario.
Estimates are bolded when significant at the alpha = 0.05 level. Each of
the 338 districts is represented in this model.}

\centering{

\centering
\begin{tblr}[         %% tabularray outer open
]                     %% tabularray outer close
{                     %% tabularray inner open
colspec={Q[]Q[]Q[]Q[]Q[]},
hline{2}={5}{solid, black, 0.03em},
hline{2}={2,4}{solid, black, 0.03em, l=-0.5},
hline{2}={3}{solid, black, 0.03em, r=-0.5},
hline{3}={1-5}{solid, black, 0.05em},
hline{17}={1-5}{solid, black, 0.05em},
hline{1}={1-5}{solid, black, 0.08em},
hline{22}={1-5}{solid, black, 0.08em},
cell{1}{1}={}{halign=l},
cell{1}{2}={c=2}{halign=c},
cell{1}{3}={}{halign=c},
cell{1}{4}={c=2}{halign=c},
cell{1}{5}={}{halign=c},
cell{2-3,6,9-10,12-14,16}{2}={}{font=\bfseries\fontsize{0.85em}{1.15em}\selectfont, halign=c},
cell{2-3,6,9-10,12-14,16}{3}={}{font=\bfseries\fontsize{0.85em}{1.15em}\selectfont, halign=c},
cell{2-3,6,9-10,12-14,16}{4}={}{font=\bfseries\fontsize{0.85em}{1.15em}\selectfont, halign=c},
cell{2-3,6,9-10,12-14,16}{5}={}{font=\bfseries\fontsize{0.85em}{1.15em}\selectfont, halign=c},
cell{2}{1}={}{font=\bfseries\fontsize{0.85em}{1.15em}\selectfont, halign=l},
cell{3-21}{1}={}{font=\fontsize{0.85em}{1.15em}\selectfont, halign=l},
cell{4-5,7-8,11,15,17-21}{2}={}{font=\fontsize{0.85em}{1.15em}\selectfont, halign=c},
cell{4-5,7-8,11,15,17-21}{3}={}{font=\fontsize{0.85em}{1.15em}\selectfont, halign=c},
cell{4-5,7-8,11,15,17-21}{4}={}{font=\fontsize{0.85em}{1.15em}\selectfont, halign=c},
cell{4-5,7-8,11,15,17-21}{5}={}{font=\fontsize{0.85em}{1.15em}\selectfont, halign=c},
}                     %% tabularray inner close
& Coefficient &  & Odds Ratio &  \\
& Est. & S.E. & Est. & S.E. \\
(Intercept) & \num{-3.939} & \num{0.862} & \num{0.019} & \num{0.017} \\
Amount & \num{-0.000} & \num{0.000} & \num{1.000} & \num{0.000} \\
Ontario & - &  & - &  \\
Quebec & \num{0.538} & \num{0.120} & \num{1.713} & \num{0.205} \\
British Columbia & \num{-0.182} & \num{0.104} & \num{0.834} & \num{0.087} \\
Prairies & \num{-0.205} & \num{0.117} & \num{0.815} & \num{0.095} \\
Atlantic Canada & \num{0.301} & \num{0.151} & \num{1.351} & \num{0.203} \\
Territories & \num{-1.222} & \num{0.497} & \num{0.295} & \num{0.147} \\
Age & \num{0.012} & \num{0.017} & \num{1.012} & \num{0.017} \\
Income & \num{0.000} & \num{0.000} & \num{1.000} & \num{0.000} \\
Education & \num{1.658} & \num{0.646} & \num{5.247} & \num{3.391} \\
Visible Minority & \num{1.847} & \num{0.238} & \num{6.340} & \num{1.511} \\
Margin of Victory (mean) & \num{0.004} & \num{0.003} & \num{1.004} & \num{0.003} \\
Population Density (sqrt) & \num{0.012} & \num{0.002} & \num{1.012} & \num{0.002} \\
Num.Obs. & \num{338} &  & \num{338} &  \\
R2 & \num{0.607} &  & \num{0.607} &  \\
AIC & \num{-401.2} &  & \num{-401.2} &  \\
BIC & \num{-343.8} &  & \num{-343.8} &  \\
RMSE & \num{0.14} &  & \num{0.14} &  \\
\end{tblr}

}

\end{table}%

\begin{table}

\caption{\label{tbl-results-model-individual-candidiate}Results from a
logistic regression model where the outcome is whether or not an
individual donation was sent out-of-district. Only donations sent to
electoral candidates during election periods are considered. The
reference party is the Conservative Party, the reference region is
Ontario. Estimates are bolded when significant at the alpha = 0.05
level.}

\centering{

\centering
\begin{tblr}[         %% tabularray outer open
]                     %% tabularray outer close
{                     %% tabularray inner open
colspec={Q[]Q[]Q[]Q[]Q[]},
hline{2}={5}{solid, black, 0.03em},
hline{2}={2,4}{solid, black, 0.03em, l=-0.5},
hline{2}={3}{solid, black, 0.03em, r=-0.5},
hline{3}={1-5}{solid, black, 0.05em},
hline{21}={1-5}{solid, black, 0.05em},
hline{1}={1-5}{solid, black, 0.08em},
hline{27}={1-5}{solid, black, 0.08em},
cell{1}{1}={}{halign=l},
cell{1}{2}={c=2}{halign=c},
cell{1}{3}={}{halign=c},
cell{1}{4}={c=2}{halign=c},
cell{1}{5}={}{halign=c},
cell{2-3,5,7-8,10-14,16-20}{2}={}{font=\bfseries\fontsize{0.85em}{1.15em}\selectfont, halign=c},
cell{2-3,5,7-8,10-14,16-20}{3}={}{font=\bfseries\fontsize{0.85em}{1.15em}\selectfont, halign=c},
cell{2-3,5,7-8,10-14,16-20}{4}={}{font=\bfseries\fontsize{0.85em}{1.15em}\selectfont, halign=c},
cell{2-3,5,7-8,10-14,16-20}{5}={}{font=\bfseries\fontsize{0.85em}{1.15em}\selectfont, halign=c},
cell{2}{1}={}{font=\bfseries\fontsize{0.85em}{1.15em}\selectfont, halign=l},
cell{3-26}{1}={}{font=\fontsize{0.85em}{1.15em}\selectfont, halign=l},
cell{4,6,9,15,21-26}{2}={}{font=\fontsize{0.85em}{1.15em}\selectfont, halign=c},
cell{4,6,9,15,21-26}{3}={}{font=\fontsize{0.85em}{1.15em}\selectfont, halign=c},
cell{4,6,9,15,21-26}{4}={}{font=\fontsize{0.85em}{1.15em}\selectfont, halign=c},
cell{4,6,9,15,21-26}{5}={}{font=\fontsize{0.85em}{1.15em}\selectfont, halign=c},
}                     %% tabularray inner close
& Coefficient &  & Odds Ratio &  \\
& Est. & S.E. & Est. & S.E. \\
(Intercept) & \num{-3.450} & \num{0.143} & \num{0.032} & \num{0.005} \\
Conservative & - &  & - &  \\
Liberal & \num{0.079} & \num{0.037} & \num{1.082} & \num{0.040} \\
N.D.P & \num{0.007} & \num{0.036} & \num{1.007} & \num{0.036} \\
Other & \num{0.618} & \num{0.045} & \num{1.856} & \num{0.083} \\
Amount & \num{0.000} & \num{0.000} & \num{1.000} & \num{0.000} \\
Ontario & - &  & - &  \\
Quebec & \num{0.398} & \num{0.048} & \num{1.489} & \num{0.071} \\
British Columbia & \num{-0.438} & \num{0.043} & \num{0.645} & \num{0.028} \\
Prairies & \num{-0.292} & \num{0.044} & \num{0.747} & \num{0.033} \\
Atlantic Canada & \num{-0.237} & \num{0.051} & \num{0.789} & \num{0.040} \\
Territories & \num{-3.012} & \num{0.507} & \num{0.049} & \num{0.025} \\
Age & \num{0.005} & \num{0.002} & \num{1.005} & \num{0.002} \\
Income & \num{0.000} & \num{0.000} & \num{1.000} & \num{0.000} \\
Education & \num{2.011} & \num{0.136} & \num{7.471} & \num{1.019} \\
Visible Minority & \num{2.636} & \num{0.063} & \num{13.960} & \num{0.878} \\
Margin of Victory (mean) & \num{0.005} & \num{0.001} & \num{1.005} & \num{0.001} \\
Population Density (sqrt) & \num{0.004} & \num{0.001} & \num{1.004} & \num{0.001} \\
Num.Obs. & \num{28995} &  & \num{28995} &  \\
AIC & \num{33320.5} &  & \num{33320.5} &  \\
BIC & \num{33452.9} &  & \num{33452.9} &  \\
Log.Lik. & \num{-16644.236} &  & \num{-16644.236} &  \\
F & \num{273.999} &  & \num{273.999} &  \\
RMSE & \num{0.44} &  & \num{0.44} &  \\
\end{tblr}

}

\end{table}%

\begin{table}

\caption{\label{tbl-results-model-individual-candidiate-competitive}Results
from a logistic regression model where the outcome is whether or not an
individual donation was sent out-of-district. Only donations sent to
electoral candidates during election periods are considered. The
reference party is the Conservative Party, the reference region is
Ontario, and the reference general election period is the 42nd.
Estimates are bolded when significant at the alpha = 0.05 level.}

\centering{

\centering
\begin{tblr}[         %% tabularray outer open
]                     %% tabularray outer close
{                     %% tabularray inner open
colspec={Q[]Q[]Q[]Q[]Q[]},
hline{2}={5}{solid, black, 0.03em},
hline{2}={2,4}{solid, black, 0.03em, l=-0.5},
hline{2}={3}{solid, black, 0.03em, r=-0.5},
hline{3}={1-5}{solid, black, 0.05em},
hline{24}={1-5}{solid, black, 0.05em},
hline{1}={1-5}{solid, black, 0.08em},
hline{30}={1-5}{solid, black, 0.08em},
cell{1}{1}={}{halign=l},
cell{1}{2}={c=2}{halign=c},
cell{1}{3}={}{halign=c},
cell{1}{4}={c=2}{halign=c},
cell{1}{5}={}{halign=c},
cell{2-3,7,9-11,13-17,19-23}{2}={}{font=\bfseries\fontsize{0.85em}{1.15em}\selectfont, halign=c},
cell{2-3,7,9-11,13-17,19-23}{3}={}{font=\bfseries\fontsize{0.85em}{1.15em}\selectfont, halign=c},
cell{2-3,7,9-11,13-17,19-23}{4}={}{font=\bfseries\fontsize{0.85em}{1.15em}\selectfont, halign=c},
cell{2-3,7,9-11,13-17,19-23}{5}={}{font=\bfseries\fontsize{0.85em}{1.15em}\selectfont, halign=c},
cell{2}{1}={}{font=\bfseries\fontsize{0.85em}{1.15em}\selectfont, halign=l},
cell{3-29}{1}={}{font=\fontsize{0.85em}{1.15em}\selectfont, halign=l},
cell{4-6,8,12,18,24-29}{2}={}{font=\fontsize{0.85em}{1.15em}\selectfont, halign=c},
cell{4-6,8,12,18,24-29}{3}={}{font=\fontsize{0.85em}{1.15em}\selectfont, halign=c},
cell{4-6,8,12,18,24-29}{4}={}{font=\fontsize{0.85em}{1.15em}\selectfont, halign=c},
cell{4-6,8,12,18,24-29}{5}={}{font=\fontsize{0.85em}{1.15em}\selectfont, halign=c},
}                     %% tabularray inner close
& Coefficient &  & Odds Ratio &  \\
& Est. & S.E. & Est. & S.E. \\
(Intercept) & \num{-3.339} & \num{0.143} & \num{0.035} & \num{0.005} \\
Conservative & - &  & - &  \\
Liberal & \num{0.059} & \num{0.037} & \num{1.061} & \num{0.039} \\
N.D.P & \num{-0.014} & \num{0.036} & \num{0.986} & \num{0.036} \\
Other & \num{0.638} & \num{0.045} & \num{1.893} & \num{0.085} \\
42nd general election & - &  & - &  \\
43rd general election & \num{-0.158} & \num{0.033} & \num{0.854} & \num{0.028} \\
44th general election & \num{-0.119} & \num{0.036} & \num{0.888} & \num{0.032} \\
Amount & \num{0.000} & \num{0.000} & \num{1.000} & \num{0.000} \\
Ontario & - &  & - &  \\
Quebec & \num{0.402} & \num{0.047} & \num{1.495} & \num{0.071} \\
British Columbia & \num{-0.449} & \num{0.043} & \num{0.638} & \num{0.028} \\
Prairies & \num{-0.260} & \num{0.042} & \num{0.771} & \num{0.032} \\
Atlantic Canada & \num{-0.224} & \num{0.052} & \num{0.799} & \num{0.041} \\
Territories & \num{-3.028} & \num{0.507} & \num{0.048} & \num{0.025} \\
Age & \num{0.005} & \num{0.002} & \num{1.005} & \num{0.002} \\
Income & \num{0.000} & \num{0.000} & \num{1.000} & \num{0.000} \\
Education & \num{1.989} & \num{0.136} & \num{7.307} & \num{0.995} \\
Visible Minority & \num{2.636} & \num{0.063} & \num{13.950} & \num{0.879} \\
Margin of Victory (per-cycle) & \num{0.003} & \num{0.001} & \num{1.003} & \num{0.001} \\
Population Density (sqrt) & \num{0.004} & \num{0.001} & \num{1.004} & \num{0.001} \\
Num.Obs. & \num{28995} &  & \num{28995} &  \\
AIC & \num{33304.4} &  & \num{33304.4} &  \\
BIC & \num{33453.3} &  & \num{33453.3} &  \\
Log.Lik. & \num{-16634.198} &  & \num{-16634.198} &  \\
F & \num{242.454} &  & \num{242.454} &  \\
RMSE & \num{0.44} &  & \num{0.44} &  \\
\end{tblr}

}

\end{table}%

\begin{table}

\caption{\label{tbl-results-model-individual-disaggregated}Results from
a logistic regression model where the outcome is whether or not an
individual donation was sent out-of-district. The reference entity is
district associations, the reference party is the Conservative Party,
the reference province is Ontario. Estimates are bolded when significant
at the alpha = 0.05 level.}

\centering{

\centering\centering
\fontsize{10}{12}\selectfont
\begin{tabular}[t]{lcccc}
\toprule
\multicolumn{1}{c}{ } & \multicolumn{2}{c}{Coefficient} & \multicolumn{2}{c}{Odds Ratio} \\
\cmidrule(l{3pt}r{3pt}){2-3} \cmidrule(l{3pt}r{3pt}){4-5}
\textbf{ } & \textbf{Est.} & \textbf{S.E.} & \textbf{Est.} & \textbf{S.E.}\\
\midrule
\textbf{(Intercept)} & \textbf{\num{-2.851}} & \textbf{\num{0.065}} & \textbf{\num{0.058}} & \textbf{\num{0.004}}\\
\textbf{Liberal} & \textbf{\num{0.328}} & \textbf{\num{0.014}} & \textbf{\num{1.388}} & \textbf{\num{0.020}}\\
N.D.P & \num{0.012} & \num{0.019} & \num{1.012} & \num{0.019}\\
\textbf{Other} & \textbf{\num{0.344}} & \textbf{\num{0.028}} & \textbf{\num{1.410}} & \textbf{\num{0.040}}\\
\textbf{Candidate} & \textbf{\num{-0.141}} & \textbf{\num{0.015}} & \textbf{\num{0.868}} & \textbf{\num{0.013}}\\
\textbf{Amount} & \textbf{\num{0.000}} & \textbf{\num{0.000}} & \textbf{\num{1.000}} & \textbf{\num{0.000}}\\
\textbf{Alberta} & \textbf{\num{-0.472}} & \textbf{\num{0.023}} & \textbf{\num{0.624}} & \textbf{\num{0.014}}\\
\textbf{British Columbia} & \textbf{\num{-0.429}} & \textbf{\num{0.019}} & \textbf{\num{0.651}} & \textbf{\num{0.012}}\\
\textbf{Manitoba} & \textbf{\num{-0.305}} & \textbf{\num{0.033}} & \textbf{\num{0.737}} & \textbf{\num{0.024}}\\
\textbf{New Brunswick} & \textbf{\num{-0.409}} & \textbf{\num{0.048}} & \textbf{\num{0.664}} & \textbf{\num{0.032}}\\
\textbf{Newfoundland and Labrador} & \textbf{\num{0.251}} & \textbf{\num{0.050}} & \textbf{\num{1.286}} & \textbf{\num{0.065}}\\
\textbf{Northwest Territories} & \textbf{\num{-1.665}} & \textbf{\num{0.259}} & \textbf{\num{0.189}} & \textbf{\num{0.049}}\\
\textbf{Nova Scotia} & \textbf{\num{-0.427}} & \textbf{\num{0.041}} & \textbf{\num{0.652}} & \textbf{\num{0.027}}\\
\textbf{Nunavut} & \textbf{\num{-2.306}} & \textbf{\num{0.423}} & \textbf{\num{0.100}} & \textbf{\num{0.042}}\\
\textbf{Prince Edward Island} & \textbf{\num{-0.373}} & \textbf{\num{0.070}} & \textbf{\num{0.689}} & \textbf{\num{0.048}}\\
\textbf{Quebec} & \textbf{\num{0.251}} & \textbf{\num{0.021}} & \textbf{\num{1.285}} & \textbf{\num{0.027}}\\
\textbf{Saskatchewan} & \textbf{\num{-0.277}} & \textbf{\num{0.038}} & \textbf{\num{0.758}} & \textbf{\num{0.029}}\\
\textbf{Yukon} & \textbf{\num{-2.110}} & \textbf{\num{0.193}} & \textbf{\num{0.121}} & \textbf{\num{0.023}}\\
Age & \num{-0.001} & \num{0.001} & \num{0.999} & \num{0.001}\\
\textbf{Income} & \textbf{\num{0.000}} & \textbf{\num{0.000}} & \textbf{\num{1.000}} & \textbf{\num{0.000}}\\
\textbf{Education} & \textbf{\num{1.759}} & \textbf{\num{0.064}} & \textbf{\num{5.807}} & \textbf{\num{0.369}}\\
\textbf{South Asian} & \textbf{\num{1.988}} & \textbf{\num{0.045}} & \textbf{\num{7.304}} & \textbf{\num{0.325}}\\
\textbf{Chinese} & \textbf{\num{2.411}} & \textbf{\num{0.056}} & \textbf{\num{11.144}} & \textbf{\num{0.627}}\\
\textbf{Black} & \textbf{\num{2.533}} & \textbf{\num{0.147}} & \textbf{\num{12.596}} & \textbf{\num{1.851}}\\
\textbf{Filipino} & \textbf{\num{1.627}} & \textbf{\num{0.154}} & \textbf{\num{5.090}} & \textbf{\num{0.781}}\\
\textbf{Arab} & \textbf{\num{2.192}} & \textbf{\num{0.194}} & \textbf{\num{8.952}} & \textbf{\num{1.737}}\\
\textbf{Latin American} & \textbf{\num{6.726}} & \textbf{\num{0.330}} & \textbf{\num{833.532}} & \textbf{\num{275.452}}\\
\textbf{Southeast Asian} & \textbf{\num{3.916}} & \textbf{\num{0.341}} & \textbf{\num{50.222}} & \textbf{\num{17.134}}\\
\textbf{West Asian} & \textbf{\num{1.421}} & \textbf{\num{0.226}} & \textbf{\num{4.140}} & \textbf{\num{0.936}}\\
\textbf{Korean} & \textbf{\num{2.514}} & \textbf{\num{0.393}} & \textbf{\num{12.354}} & \textbf{\num{4.860}}\\
\textbf{Japanese} & \textbf{\num{3.253}} & \textbf{\num{0.824}} & \textbf{\num{25.859}} & \textbf{\num{21.304}}\\
\textbf{Margin of Victory (mean)} & \textbf{\num{0.008}} & \textbf{\num{0.001}} & \textbf{\num{1.008}} & \textbf{\num{0.001}}\\
\textbf{Population Density (sqrt)} & \textbf{\num{0.003}} & \textbf{\num{0.000}} & \textbf{\num{1.003}} & \textbf{\num{0.000}}\\
\midrule
Num.Obs. & \num{136437} &  & \num{136437} & \\
F & \num{562.004} &  & \num{562.004} & \\
RMSE & \num{0.46} &  & \num{0.46} & \\
\bottomrule
\end{tabular}

}

\end{table}%

\begin{table}

\caption{\label{tbl-results-model-district-disaggregated}Results from a
beta regression model where the outcome is the out-of-district share of
the sum of contributions over \$200 originating from an individual
district during our study period. The reference province is Ontario.
Estimates are bolded when significant at the alpha = 0.05 level. Each of
the 338 districts is represented in this model.}

\centering{

\centering\centering
\fontsize{10}{12}\selectfont
\begin{tabular}[t]{lcccc}
\toprule
\multicolumn{1}{c}{ } & \multicolumn{2}{c}{Coefficient} & \multicolumn{2}{c}{Odds Ratio} \\
\cmidrule(l{3pt}r{3pt}){2-3} \cmidrule(l{3pt}r{3pt}){4-5}
\textbf{ } & \textbf{Est.} & \textbf{S.E.} & \textbf{Est.} & \textbf{S.E.}\\
\midrule
\textbf{(Intercept)} & \textbf{\num{-4.926}} & \textbf{\num{0.875}} & \textbf{\num{0.007}} & \textbf{\num{0.006}}\\
Amount & \num{-0.000} & \num{0.000} & \num{1.000} & \num{0.000}\\
Alberta & \num{-0.238} & \num{0.134} & \num{0.788} & \num{0.105}\\
British Columbia & \num{-0.061} & \num{0.140} & \num{0.941} & \num{0.132}\\
Manitoba & \num{0.141} & \num{0.180} & \num{1.152} & \num{0.207}\\
New Brunswick & \num{0.307} & \num{0.210} & \num{1.360} & \num{0.286}\\
\textbf{Newfoundland and Labrador} & \textbf{\num{0.518}} & \textbf{\num{0.240}} & \textbf{\num{1.679}} & \textbf{\num{0.403}}\\
Northwest Territories & \num{-1.134} & \num{0.774} & \num{0.322} & \num{0.249}\\
Nova Scotia & \num{0.157} & \num{0.199} & \num{1.170} & \num{0.232}\\
Nunavut & \num{-0.707} & \num{0.972} & \num{0.493} & \num{0.479}\\
Prince Edward Island & \num{-0.038} & \num{0.318} & \num{0.963} & \num{0.306}\\
\textbf{Quebec} & \textbf{\num{0.483}} & \textbf{\num{0.118}} & \textbf{\num{1.622}} & \textbf{\num{0.191}}\\
Saskatchewan & \num{0.171} & \num{0.189} & \num{1.187} & \num{0.224}\\
\textbf{Yukon} & \textbf{\num{-1.633}} & \textbf{\num{0.775}} & \textbf{\num{0.195}} & \textbf{\num{0.151}}\\
Age & \num{0.023} & \num{0.017} & \num{1.023} & \num{0.017}\\
\textbf{Income} & \textbf{\num{0.000}} & \textbf{\num{0.000}} & \textbf{\num{1.000}} & \textbf{\num{0.000}}\\
\textbf{Education} & \textbf{\num{1.977}} & \textbf{\num{0.722}} & \textbf{\num{7.218}} & \textbf{\num{5.212}}\\
\textbf{South Asian} & \textbf{\num{1.562}} & \textbf{\num{0.411}} & \textbf{\num{4.771}} & \textbf{\num{1.961}}\\
\textbf{Chinese} & \textbf{\num{2.250}} & \textbf{\num{0.444}} & \textbf{\num{9.484}} & \textbf{\num{4.215}}\\
\textbf{Black} & \textbf{\num{4.204}} & \textbf{\num{1.206}} & \textbf{\num{66.976}} & \textbf{\num{80.740}}\\
Filipino & \num{0.653} & \num{1.191} & \num{1.921} & \num{2.287}\\
Arab & \num{2.687} & \num{1.567} & \num{14.687} & \num{23.011}\\
Latin American & \num{3.856} & \num{4.289} & \num{47.256} & \num{202.694}\\
Southeast Asian & \num{6.267} & \num{4.493} & \num{527.095} & \num{2368.115}\\
West Asian & \num{0.402} & \num{2.303} & \num{1.494} & \num{3.442}\\
Korean & \num{4.480} & \num{4.205} & \num{88.201} & \num{370.871}\\
Japanese & \num{-6.419} & \num{17.602} & \num{0.002} & \num{0.029}\\
Margin of Victory (mean) & \num{0.004} & \num{0.003} & \num{1.004} & \num{0.003}\\
\textbf{Population Density (sqrt)} & \textbf{\num{0.010}} & \textbf{\num{0.003}} & \textbf{\num{1.010}} & \textbf{\num{0.003}}\\
\midrule
Num.Obs. & \num{338} &  & \num{338} & \\
R2 & \num{0.683} &  & \num{0.683} & \\
RMSE & \num{0.12} &  & \num{0.12} & \\
\bottomrule
\end{tabular}

}

\end{table}%

\section{Testing Whether Out-of-District Rates can be Explained by
Distance Alone}\label{sec-district-propensity-simulation}

\subsection{Simulation Setup}\label{simulation-setup}

Figure~\ref{fig-data-summaries-ood-rate-district-scatterplot} shows how
variable out-of-district donation rates are at the district level.
However, this may simply be a byproduct of the spatial arrangement of
electoral districts. Districts vary enormously in size, from the 13
square kilometres that make up Toronto--St.~Paul's, to the entire
territory of Nunavut. Assuming that social influence decays as a
function of distance, we should expect a donor residing in a district in
downtown Toronto to (1) be more aware of or (2) perceive to have a
greater stake in district associations and candidates outside their own
district, simply by virtue of their spatial proximity. By contrast, a
donor in Nunavut may be expected to place less importance on political
entities in the neighbouring Northwest Territories, whose population
centre is nearly one thousand kilometres away. In this section we design
a simulation to test the null hypothesis that the observed variance in
district-level out-of-district donation rates is well-explained by
differences in spatial clustering.

\begin{itemize}
\tightlist
\item
  Assign a random distance to each donation.
\item
  If this distance is \textless{} 1km, immediately declare the donation
  as having been received in-district. Otherwise, continue.
\item
  Assign a random direction for that donation.
\item
  Determine if it landed within Canada. If it didn't choose a different
  direction. Otherwise, continue.
\item
  Determine which district it landed in.
\item
  Declare in-district vs.~out-of-district.
\end{itemize}

Distances were assigned by sampling from the true distance distribution
of the observed donations dataset, shown in
Figure~\ref{fig-district-propensity-simulation-distance-distribution}
(we note that this figure only shows out-of-district donations; the
majority of donations were in-district, and therefore travelled zero
kilometres). The distance a donation travelled was defined to be zero if
the donation was in-district, and otherwise set to be the spatial
distance between the centroid of the donor's dissemination area, and the
centroid of the recipient's district. Shapefiles were used to determine
which district a randomly travelling donation landed in.

For each district, 30 replicate simulations were run to estimate the
sample distribution of the simulation. Using these distributions, the
observed out-of-district proportions were compared with the simulated
out-of-district proportions, for each district. Only results for the
count of out-of-district donors (not the amount of donations sent
out-of-district) was considered, however results were very similar for
amounts, just with systematically wider sampling distributions.

\begin{figure}

\centering{

\pandocbounded{\includegraphics[keepaspectratio]{paper_files/figure-pdf/fig-district-propensity-simulation-distance-distribution-1.pdf}}

}

\caption{\label{fig-district-propensity-simulation-distance-distribution}Distribution
of donation distances for out-of-district donations over \$200 during
our study period. Both donations to candidates and district associations
are shown together, however distributions were very similar for both
entities separately. Distance is measured as the distance from the donor
dissemination area's centroid to the recipient district's centroid.}

\end{figure}%

\subsection{Simulation Results}\label{simulation-results}

When considering donations addressed to candidates (aggregated over all
three election periods), only 6.5\% of districts fell within the
expected proportion of out-of-district donations. For donations
addressed to EDAs over the entire study period, only 5.3\% of districts
fell within the expected proportion of out-of-district donations.
Proportions of districts that fell within their expected ranges were
slightly higher (8.2\% and 8.5\% respectively) when looking at the share
of the total contributed amount, however this was likely due to the
systematically wider sampling distributions caused by the added variance
incurred when using amounts.

Figure~\ref{fig-district-propensity-simulation-distance-deviations}
shows that, for both entities, most districts have out-of-district
proportions that deviate drastically from simulated expectations under
the null hypothesis. Further, most deviations, including the largest
ones in magnitude, were districts where the observed out-of-district
donation rate was larger than expected, as opposed to smaller. This
indicates that most districts have higher out-of-district donation rates
than what can be explained by the spatial arrangement of districts
alone. This can be seen more clearly in
Figure~\ref{fig-district-propensity-simulation-distribution-comparison};
in both cases, much of the support for the observed distribution lies
strictly above the support for the expected distribution. This skew in
deviations is especially noticeable for donations addressed to district
associations. Results using amount-weighted proportions were similar in
both cases.

\begin{figure}

\centering{

\pandocbounded{\includegraphics[keepaspectratio]{paper_files/figure-pdf/fig-district-propensity-simulation-distance-deviations-1.pdf}}

}

\caption{\label{fig-district-propensity-simulation-distance-deviations}Deviation
of each district's observed proportion of out-of-district (OOD)
donations from its simulated proportion under the null hypothesis.
Results for both donations to candidates (A) and district associations
(B) are shown. Districts are sorted by the sign and magnitude of their
deviation. The dashed red line indicates no deviation, and the shaded
red band gives the simulated 95\% confidence interval. Points outside
the band deviate more than spatial clustering alone can explain.}

\end{figure}%

\begin{figure}

\centering{

\pandocbounded{\includegraphics[keepaspectratio]{paper_files/figure-pdf/fig-district-propensity-simulation-distribution-comparison-1.pdf}}

}

\caption{\label{fig-district-propensity-simulation-distribution-comparison}Distribution
across districts of the observed (black) and simulated null-expected
(red) proportion of out-of-district donations. Both donations to
candidates (A) and district associations (B) are shown.}

\end{figure}%

\section{Additional Data Summaries}\label{sec-additional-data-summaries}

\begin{table}

\caption{\label{tbl-additional-data-summaries-situation}Summaries of
financial contributions to federal political entities over \$200 from
2015 to 2024, separated by entity. Only donations stated to have been
received during a general election period are shown. Contributions
received by nomination contestants are not shown.}

\centering{

\centering
\begin{tblr}[         %% tabularray outer open
]                     %% tabularray outer close
{                     %% tabularray inner open
width={1\linewidth},
colspec={X[]X[]X[]X[]X[]},
hline{2}={1-5}{solid, black, 0.05em},
hline{1}={1-5}{solid, black, 0.08em},
hline{6}={1-5}{solid, black, 0.08em},
cell{1}{1}={}{font=\bfseries\fontsize{0.85em}{1.15em}\selectfont, halign=l},
cell{1}{2}={}{font=\bfseries\fontsize{0.85em}{1.15em}\selectfont, halign=l},
cell{1}{3}={}{font=\bfseries\fontsize{0.85em}{1.15em}\selectfont, halign=r},
cell{1}{4}={}{font=\bfseries\fontsize{0.85em}{1.15em}\selectfont, halign=r},
cell{1}{5}={}{font=\bfseries\fontsize{0.85em}{1.15em}\selectfont, halign=r},
cell{2-5}{1}={}{font=\fontsize{0.85em}{1.15em}\selectfont, halign=l},
cell{2-5}{2}={}{font=\fontsize{0.85em}{1.15em}\selectfont, halign=l},
cell{2-5}{3}={}{font=\fontsize{0.85em}{1.15em}\selectfont, halign=r},
cell{2-5}{4}={}{font=\fontsize{0.85em}{1.15em}\selectfont, halign=r},
cell{2-5}{5}={}{font=\fontsize{0.85em}{1.15em}\selectfont, halign=r},
}                     %% tabularray inner close
Political Entity & Type & Donation Count & Donation Amount (Thousands) & Share of Total Amount \\
Candidate & Localized & 29,380 & \$18,872 & 25.7\% \\
District Association & Localized & 36,489 & \$20,373 & 27.7\% \\
Leadership Contestant & National & 134 & \$103 & 0.1\% \\
Party & National & 79,052 & \$34,071 & 46.4\% \\
\end{tblr}

}

\end{table}%

\begin{figure}

\centering{

\pandocbounded{\includegraphics[keepaspectratio]{paper_files/figure-pdf/fig-local-vs-national-province-1.pdf}}

}

\caption{\label{fig-local-vs-national-province}Distribution of
contributions across donor provinces, shown separately for localized
entities (candidates and district associations) and national ones
(parties and leadership contestants). Within each group, bars give each
province's share of that group's total contributed amount, and sum to
100\% across provinces. Contributions over \$200 received from 2015 to
2024.}

\end{figure}%

\begin{figure}

\centering{

\pandocbounded{\includegraphics[keepaspectratio]{paper_files/figure-pdf/fig-local-vs-national-donor-district-1.pdf}}

}

\caption{\label{fig-local-vs-national-donor-district}Total amount
contributed to localized entities (candidates and district associations)
against the total contributed to national entities (parties and
leadership contestants), with one point per donor district.
Contributions over \$200 received from 2015 to 2024.}

\end{figure}%

\begin{figure}

\centering{

\pandocbounded{\includegraphics[keepaspectratio]{paper_files/figure-pdf/fig-ood-rate-amount-distributions-1.pdf}}

}

\caption{\label{fig-ood-rate-amount-distributions}Value distributions of
in- (blue) and out-of-district (red) contributions over \$200 received
by candidates (left) and district associations (right) from 2015 to
2024. Dashed grey lines mark notable individual contribution limits.}

\end{figure}%

\begin{table}

\caption{\label{tbl-data-summaries-ood-rate-party}Proportion of
contributions over \$200 that were received by localized entities
(candidates and district associations) outside the donor's district,
from 2015 to 2024, broken down by recipient entity and party. The three
largest parties are shown individually; all others are grouped into
`Other'. Raw and value-weighted proportions of out-of-district donations
(OOD) are reported, alongside overall contribution counts for reference.
Contributions received by candidates outside general election cycles are
not included in calculations.}

\centering{

\centering
\begin{tblr}[         %% tabularray outer open
]                     %% tabularray outer close
{                     %% tabularray inner open
width={1\linewidth},
colspec={X[]X[]X[]X[]X[]},
hline{2}={1-5}{solid, black, 0.05em},
hline{1}={1-5}{solid, black, 0.08em},
hline{10}={1-5}{solid, black, 0.08em},
cell{1}{1}={}{font=\bfseries\fontsize{0.85em}{1.15em}\selectfont, halign=l},
cell{1}{2}={}{font=\bfseries\fontsize{0.85em}{1.15em}\selectfont, halign=l},
cell{1}{3}={}{font=\bfseries\fontsize{0.85em}{1.15em}\selectfont, halign=r},
cell{1}{4}={}{font=\bfseries\fontsize{0.85em}{1.15em}\selectfont, halign=r},
cell{1}{5}={}{font=\bfseries\fontsize{0.85em}{1.15em}\selectfont, halign=r},
cell{2-9}{1}={}{font=\fontsize{0.85em}{1.15em}\selectfont, halign=l},
cell{2-9}{2}={}{font=\fontsize{0.85em}{1.15em}\selectfont, halign=l},
cell{2-9}{3}={}{font=\fontsize{0.85em}{1.15em}\selectfont, halign=r},
cell{2-9}{4}={}{font=\fontsize{0.85em}{1.15em}\selectfont, halign=r},
cell{2-9}{5}={}{font=\fontsize{0.85em}{1.15em}\selectfont, halign=r},
}                     %% tabularray inner close
Entity & Party & Donation Count & OOD Proportion & OOD Proportion (Amount Weighted) \\
Candidate & Conservative Party of Canada & 14,421 & 36.1\% & 38.5\% \\
Candidate & Liberal Party of Canada & 5,961 & 39.7\% & 43.9\% \\
Candidate & New Democratic Party & 6,637 & 35.9\% & 39.6\% \\
Candidate & Other & 3,540 & 43.3\% & 41.4\% \\
District Association & Conservative Party of Canada & 38,756 & 36.0\% & 40.1\% \\
District Association & Liberal Party of Canada & 49,859 & 53.5\% & 57.8\% \\
District Association & New Democratic Party & 16,164 & 35.2\% & 35.8\% \\
District Association & Other & 3,840 & 34.6\% & 36.0\% \\
\end{tblr}

}

\end{table}%

\begin{figure}

\centering{

\pandocbounded{\includegraphics[keepaspectratio]{paper_files/figure-pdf/fig-data-summaries-ood-rate-district-distribution-1.pdf}}

}

\caption{\label{fig-data-summaries-ood-rate-district-distribution}Distribution
of out-of-district (OOD) contribution shares for each district in our
study (for contributions over \$200, from 2015 to 2024). (A) and (B)
show this distribution for donations to district associations and
candidates separately. The district-level mean is indicated with an
orange vertical line in each case. Distributions were similar when raw
instead of amount-weighted proportions were used.}

\end{figure}%

\begin{figure}

\centering{

\pandocbounded{\includegraphics[keepaspectratio]{paper_files/figure-pdf/fig-data-summaries-ood-rate-province-distribution-1.pdf}}

}

\caption{\label{fig-data-summaries-ood-rate-province-distribution}Out-of-district
share of the total amount contributed over \$200 to (A) district
associations and (B) candidates, for each donor province, from 2015 to
2024. Within each panel, provinces are ordered from highest to lowest
out-of-district share.}

\end{figure}%

\begin{figure}

\centering{

\pandocbounded{\includegraphics[keepaspectratio]{paper_files/figure-pdf/fig-data-summaries-ood-rate-writ-periods-scatterplot-1.pdf}}

}

\caption{\label{fig-data-summaries-ood-rate-writ-periods-scatterplot}Relationship
between the sum of contributions over \$200 donated in- v.s.
out-of-district, for each district during the 2013 representational
order (effective between 2015 and 2024). (A), (B), and (C) display this
relationship for donations to candidates during the 42nd, 43rd, and 44th
writ periods respectively. Points represent individual districts. The
grey dashed line is 1:1. The top 3 districts by amount contributed in-
and out-of-district respectively are labelled with their names.}

\end{figure}%

\newpage

\section*{References}\label{references}
\addcontentsline{toc}{section}{References}

\protect\phantomsection\label{refs}
\begin{CSLReferences}{1}{1}
\bibitem[\citeproctext]{ref-ansolabehere_why_2003}
Ansolabehere, Stephen, John M. de Figueiredo, and James M. Snyder Jr.
2003. {``Why Is {There} so {Little} {Money} in {U}.{S}. {Politics}?''}
\emph{Journal of Economic Perspectives} 17 (1): 105--30.
\url{https://doi.org/10.1257/089533003321164976}.

\bibitem[\citeproctext]{ref-baker_getting_2016}
Baker, Anne E. 2016. {``Getting {Short}-{Changed}? {The} {Impact} of
{Outside} {Money} on {District} {Representation}.''} \emph{Social
Science Quarterly} 97 (5): 1096--107.
\url{https://www.jstor.org/stable/26612374}.

\bibitem[\citeproctext]{ref-bednar_political_2011}
Bednar, Jenna, and Elizabeth R. Gerber. 2011. \emph{Political
{Geography}, {Campaign} {Contributions}, and {Representation}}.
\url{https://public.websites.umich.edu/~jbednar/WIP/irod_050211_wfigs.pdf}.

\bibitem[\citeproctext]{ref-besco_identities_2019}
Besco, Randy. 2019. \emph{Identities and {Interests}: {Race},
{Ethnicity}, and {Affinity} {Voting}}. UBC Press.
\url{https://www.ubcpress.ca/identities-and-interests}.

\bibitem[\citeproctext]{ref-besco_ethnic_2022}
Besco, Randy, and Erin Tolley. 2022. {``Ethnic Group Differences in
Donations to Electoral Candidates.''} \emph{Journal of Ethnic and
Migration Studies} 48 (5): 1072--94.
\url{https://doi.org/10.1080/1369183X.2020.1804339}.

\bibitem[\citeproctext]{ref-bramlett_political_2011}
Bramlett, Brittany H., James G. Gimpel, and Frances E. Lee. 2011. {``The
{Political} {Ecology} of {Opinion} in {Big}-{Donor} {Neighborhoods}.''}
\emph{Political Behavior} 33 (4): 565--600.
\url{https://doi.org/10.1007/s11109-010-9144-7}.

\bibitem[\citeproctext]{ref-canada_post_guidelines_2026}
Canada Post. 2026. \emph{Addressing Guidelines: Overview}.
\url{https://www.canadapost-postescanada.ca/cpc/en/support/articles/addressing-guidelines/overview.page\#1414221}.

\bibitem[\citeproctext]{ref-canes_wrone_out_2022}
Canes-Wrone, Brandice, and Kenneth M. Miller. 2022. {``Out-of-{District}
{Contributors} and {Representation} in the {US} {House}.''}
\emph{Legislative Studies Quarterly}.
\url{https://csap.yale.edu/sites/default/files/files/apppw-bcwrevised-2-5-20.pdf}.

\bibitem[\citeproctext]{ref-carmichael_political_2014}
Carmichael, Brianna, and Paul Howe. 2014. {``Political {Donations} and
{Democratic} {Equality} in {Canada}.''} \emph{Canadian Parliamentary
Review}. \url{http://www.revparl.ca/37/1/37n1e_14_carmichael.pdf}.

\bibitem[\citeproctext]{ref-corrado_new_2010}
Corrado, Anthony, Thomas E. Mann, Daniel R. Ortiz, and Trevor Potter.
2010. \emph{The New Campaign Finance Sourcebook}. Brookings Institution
Press.

\bibitem[\citeproctext]{ref-betareg}
Cribari-Neto, Francisco, and Achim Zeileis. 2010. {``Beta Regression in
r.''} \emph{Journal of Statistical Software} 34 (2).
\url{https://doi.org/10.18637/jss.v034.i02}.

\bibitem[\citeproctext]{ref-eagles_effectiveness_2004}
Eagles, Munroe. 2004. {``The {Effectiveness} of {Local} {Campaign}
{Spending} in the 1993 and 1997 {Federal} {Elections} in {Canada}.''}
\emph{Canadian Journal of Political Science / Revue Canadienne de
Science Politique} 37 (1): 117--36.
\url{https://www.jstor.org/stable/25165603}.

\bibitem[\citeproctext]{ref-canada_redistribution_effect_2023}
Elections Canada. 2023. \emph{Implementation of New Federal Electoral
Boundaries}. Government of Canada.
\url{https://www.elections.ca/content.aspx?section=med&dir=pre&document=sep2723&lang=e}.

\bibitem[\citeproctext]{ref-canada_open_2025}
Elections Canada. 2025a. \emph{Open {Data}}.
\url{https://www.elections.ca/content.aspx?section=fin&dir=oda&document=index&lang=e}.

\bibitem[\citeproctext]{ref-canada_redistribution_2025}
Elections Canada. 2025b. \emph{Redistribution of {Federal} {Electoral}
{Districts} 2022}.
\url{https://www.elections.ca/content.aspx?section=res&dir=cir/red&document=index&lang=e}.

\bibitem[\citeproctext]{ref-canada_annual_2026}
Elections Canada. 2026a. \emph{Annual {Limits} on {Contributions} --
2026}.
\url{https://www.elections.ca/content.aspx?section=pol&dir=lim&document=lim2026&lang=e}.

\bibitem[\citeproctext]{ref-canada_expenses_2026}
Elections Canada. 2026b. \emph{Expenses {Limits}}.
\url{https://www.elections.ca/content.aspx?section=pol&dir=limits&document=index&lang=e}.

\bibitem[\citeproctext]{ref-canada_audit_2026}
Elections Canada. 2026c. \emph{Financial Returns}.
\url{https://www.elections.ca/content.aspx?section=fin&dir=rev&document=index&lang=e}.

\bibitem[\citeproctext]{ref-canada_political_2026}
Elections Canada. 2026d. \emph{Political {Entities}}.
\url{https://www.elections.ca/content.aspx?section=pol&document=index&lang=e}.

\bibitem[\citeproctext]{ref-canada_financing_2026}
Elections Canada. 2026e. \emph{Political Financing}.
\url{https://www.elections.ca/content.aspx?section=fin&document=index&lang=e}.

\bibitem[\citeproctext]{ref-canada_understanding_2026}
Elections Canada. 2026f. \emph{Understanding Contributions}.
\url{https://www.elections.ca/content.aspx?section=fin&dir=cont&document=index&lang=e}.

\bibitem[\citeproctext]{ref-fec_contribution_limits_2026}
Federal Election Commission. 2026. {``Contribution Limits.''}
\url{https://www.fec.gov/help-candidates-and-committees/candidate-taking-receipts/contribution-limits/}.

\bibitem[\citeproctext]{ref-garnett_where_2024}
Garnett, Holly Ann. 2024. {``Where Do Donors Come from? {Using} Census
Data to Predict Donations to {Canadian} Federal Election Candidates.''}
\emph{Political Geography} 110 (April): 103077.
\url{https://doi.org/10.1016/j.polgeo.2024.103077}.

\bibitem[\citeproctext]{ref-garnett_lifeblood_2022}
Garnett, Holly Ann, Scott Pruysers, Lisa Young, and William P. Cross.
2022. {``Lifeblood of the {Party}: {Motivations} for {Political}
{Donations} in {Canada}.''} \emph{American Review of Canadian Studies}
52 (4). \url{https://doi.org/10.1080/02722011.2022.2147756}.

\bibitem[\citeproctext]{ref-gimpel_political_2006}
Gimpel, James G., Frances E. Lee, and Joshua Kaminski. 2006. {``The
{Political} {Geography} of {Campaign} {Contributions} in {American}
{Politics}.''} \emph{The Journal of Politics} 68 (3): 626--39.
\url{https://doi.org/10.1111/j.1468-2508.2006.00450.x}.

\bibitem[\citeproctext]{ref-gimpel_check_2008}
Gimpel, James G., Frances E. Lee, and Shanna Pearson-Merkowitz. 2008.
{``The {Check} {Is} in the {Mail}: {Interdistrict} {Funding} {Flows} in
{Congressional} {Elections}.''} \emph{American Journal of Political
Science} 52 (2): 373--94. \url{https://www.jstor.org/stable/25193819}.

\bibitem[\citeproctext]{ref-hou_recent_2004}
Hou, Feng. 2004. \emph{Recent {Immigration} and the {Formation} of
{Visible} {Minority} {Neighbourhoods} in {Canada}'s {Large} {Cities}}.
Analytical Studies Branch Research Paper Series 11F0019MIE No. 221.
Statistics Canada.
\url{https://www150.statcan.gc.ca/n1/pub/11f0019m/11f0019m2004221-eng.pdf}.

\bibitem[\citeproctext]{ref-jacobs_nationalization_2024}
Jacobs, Nicholas, and Wasike Gil Imboywa. 2024. {``The {Nationalization}
of {Individual} {Campaign} {Contributions} in {U}.{S}. {Senate}
{Elections}, 1984-2020.''} \emph{American Politics Research} 52 (3):
239--48. \url{https://doi.org/10.1177/1532673X231220639}.

\bibitem[\citeproctext]{ref-jansen_donates_2012}
Jansen, Harold J., Melanee Thomas, and Lisa Young. 2012. {``Who Donates
to Canada's Political Parties?''} \emph{Annual Meeting of the Canadian
Political Science Association} (Edmonton, Alberta), June.
\url{https://www.cpsa-acsp.ca/papers-2012/Jansen-Thomas-Young.pdf}.

\bibitem[\citeproctext]{ref-keena_out_2024}
Keena, Alex. 2024. {``Out-{Of}-{State} {Donors} and {Legislative}
{Surrogacy} in the {U}.{S}. {Senate}.''} \emph{Political Research
Quarterly} 77 (3): 880--90.
\url{https://doi.org/10.1177/10659129241249171}.

\bibitem[\citeproctext]{ref-lee_ethnoracial_2017}
Lee, Barrett A., and Gregory Sharp. 2017. {``Ethnoracial {Diversity}
Across the {Rural}-{Urban} {Continuum}.''} \emph{The ANNALS of the
American Academy of Political and Social Science} 672 (1): 26--45.
\url{https://doi.org/10.1177/0002716217708560}.

\bibitem[\citeproctext]{ref-mcmahon_covid_2023}
McMahon, Nicole, Anthony Sayers, and Christopher Alcantara. 2023.
{``Political Donations and the Gender Gap During {COVID}-19.''}
\emph{Party Politics} 29 (1): 176--84.
\url{https://doi.org/10.1177/13540688211047768}.

\bibitem[\citeproctext]{ref-mcmahon_follow_2015}
McMahon, Tamsin, and Julia Wolfe. 2015. {``Follow the Money: {Mapping}
{Canada}'s Richest Regions for Political Donations.''} In \emph{The
Globe and Mail}.
\url{https://www.theglobeandmail.com/news/politics/follow-the-money-mapping-canadas-richest-regions-for-political-donations/article26822168/}.

\bibitem[\citeproctext]{ref-nathan_donor_2025}
Nathan, Charles, Arvind Krishnamurthy, Curtis Bram, and Jason Douglas
Todd. 2025. {``The {Donor} {Went} {Down} to {Georgia}:
{Out}-of-{District} {Donations} and {Rivalrous} {Representation}.''}
\emph{Political Behavior} 47 (1): 41--60.
\url{https://doi.org/10.1007/s11109-024-09940-y}.

\bibitem[\citeproctext]{ref-peoples_impact_2008}
Peoples, Clayton D., and Michael Gortari. 2008. {``The Impact of
Campaign Contributions on Policymaking in {The} {U}.{S}. And {Canada}:
{Theoretical} and Public Policy Implications.''} In \emph{Politics and
{Public} {Policy}}, edited by Harland Prechel, vol. 17. Emerald Group
Publishing Limited. \url{https://doi.org/10.1016/S0895-9935(08)17003-6}.

\bibitem[\citeproctext]{ref-pruysers_leadership_2024}
Pruysers, Scott. 2024. {``Sour Grapes? Party Donors and Canadian
Leadership Primaries.''} \emph{Government and Opposition} 59 (2):
566--81. \url{https://doi.org/10.1017/gov.2023.15}.

\bibitem[\citeproctext]{ref-quillian_segregation_2012}
Quillian, Lincoln. 2012. {``Segregation and {Poverty} {Concentration}:
{The} {Role} of {Three} {Segregations}.''} \emph{American Sociological
Review} 77 (3): 354--79. \url{https://doi.org/10.1177/0003122412447793}.

\bibitem[\citeproctext]{ref-r_citation}
R Core Team. 2025. \emph{R: A Language and Environment for Statistical
Computing}. R Foundation for Statistical Computing.
\url{https://www.R-project.org/}.

\bibitem[\citeproctext]{ref-sievert_out_2023}
Sievert, Joel, and Stephanie Mathiasen. 2023. {``Out-of-{State} {Donors}
and {Nationalized} {Politics} in {U}.{S}. {Senate} {Elections}.''}
\emph{The Forum} 21 (2): 309--28.
\url{https://doi.org/10.1515/for-2023-2018}.

\bibitem[\citeproctext]{ref-generalized_beta_regression}
Simas, Alexandre B., Wagner Barreto-Souza, and Andréa V. Rocha. 2010.
{``Improved Estimators for a General Class of Beta Regression Models.''}
\emph{Computational Statistics \& Amp; Data Analysis} 54 (2): 348--66.
\url{https://doi.org/10.1016/j.csda.2009.08.017}.

\bibitem[\citeproctext]{ref-statcan_da_2021}
Statistics Canada. 2021. \emph{Dissemination Area ({DA})}. Dictionary,
Census of Population, 2021.
\url{https://www12.statcan.gc.ca/census-recensement/2021/ref/dict/az/definition-eng.cfm?ID=geo021}.

\bibitem[\citeproctext]{ref-statcan_boundary_2021}
Statistics Canada. 2022a. \emph{Boundary Files, 2021 Census}. Catalogue
no. 92-160-X. Government of Canada.
\url{https://www12.statcan.gc.ca/census-recensement/2021/geo/sip-pis/boundary-limites/index2021-eng.cfm?year=21}.

\bibitem[\citeproctext]{ref-statcan_census_zip_2022}
Statistics Canada. 2022b. \emph{Census Profile, 2021 {Census} of
Population -- Download Files}.
\url{https://www12.statcan.gc.ca/census-recensement/2021/dp-pd/prof/details/download-telecharger.cfm?Lang=E}.

\bibitem[\citeproctext]{ref-statcan_pccf_2025}
Statistics Canada. 2025. \emph{Postal Code Conversion File, June 2024,
Census of Canada 2021}. Version V2. Borealis.
\url{https://doi.org/10.5683/SP3/IT8RET}.

\bibitem[\citeproctext]{ref-tolley_who_2022}
Tolley, Erin, Randy Besco, and Semra Sevi. 2022. {``Who {Controls} the
{Purse} {Strings}? {A} {Longitudinal} {Study} of {Gender} and
{Donations} in {Canadian} {Politics}.''} \emph{Politics \& Gender} 18
(1): 244--72. \url{https://doi.org/10.1017/S1743923X20000276}.

\bibitem[\citeproctext]{ref-canada_open_government_2026}
Treasury Board of Canada Secretariat. 2026. \emph{Open {Government}}.
\url{https://open.canada.ca/en}.

\bibitem[\citeproctext]{ref-walks_ghettos_2006}
Walks, R. Alan, and Larry S. Bourne. 2006. {``Ghettos in {Canada}'s
Cities? {Racial} Segregation, Ethnic Enclaves and Poverty Concentration
in {Canadian} Urban Areas.''} \emph{Canadian Geographies / Géographies
Canadiennes} 50 (3): 273--97.
\url{https://doi.org/10.1111/j.1541-0064.2006.00142.x}.

\bibitem[\citeproctext]{ref-young_money_2011}
Young, Lisa, and Harold J. Jansen, eds. 2011. \emph{Money, {Politics},
and {Democracy}: {Canada}'s {Party} {Finance} {Reforms}}. University of
British Columbia Press. \url{https://doi.org/10.59962/9780774818933}.

\end{CSLReferences}




\end{document}
